\documentclass{article}
% \usepackage[notextcomp,nomath]{stix2}
\usepackage{tcolorbox}
\usepackage{amsfonts,amscd,amssymb,amsmath,amsthm,enumerate}
\numberwithin{equation}{section}
\usepackage{array,mathtools,mathrsfs,bm}
\usepackage{threeparttable}
\usepackage{booktabs,caption,longtable,multicol,multirow}
\usepackage{subcaption,lscape}
\usepackage{makecell}
\usepackage{graphicx}
\usepackage{enumerate}
\usepackage{enumitem}
\setlist[enumerate]{noitemsep, topsep=2pt}
\setlist[itemize]{noitemsep, topsep=2pt}
% \usepackage[linesnumbered,boxruled,lined,commentsnumbered]{algorithm2e}
\usepackage[colorlinks,citecolor=purple,linkcolor=blue]{hyperref}
% \usepackage[nameinlink]{cleveref}
\usepackage{cleveref}
% \usepackage[parfill]{parskip}
\usepackage[margin=1.5in]{geometry}
\usepackage[sort&compress,numbers]{natbib}
% \usepackage{authblk}
\usepackage{algorithm,algorithmic}
\usepackage{pifont}
\usepackage{appendix}
\usepackage{multirow}
\usepackage{tabularx}
\usepackage{makecell}
\usepackage{tikz}
\usepackage{pgfplots}
\usetikzlibrary{arrows}
\usetikzlibrary{calc}
\usetikzlibrary{arrows.meta}
\usetikzlibrary{backgrounds}
\usepgfplotslibrary{patchplots}
\usepgfplotslibrary{fillbetween}
% \usepackage[inline]{showlabels}
% \renewcommand{\showlabelfont}{\tiny\color{blue}}
%%%%%%%%%%%%%%%%%%%%%%%%%%%%%%%%%%%%%%%%%%%%%%%%
\usepackage{eqparbox}
\renewcommand\algorithmiccomment[1]{%
  \hfill\texttt{\#}\ \eqparbox{COMMENT}{\texttt{#1}}%
}
%%%%%%%%%%%%%%%%%%%%%%%%%%%%%%%%%%%%%%%%%%%%%%%%

%% tikz
\definecolor{ao(english)}{rgb}{0.0, 0.5, 0.0}
\definecolor{cadmiumgreen}{rgb}{0.0, 0.42, 0.24}
\definecolor{darkpastelgreen}{rgb}{0.01, 0.75, 0.24}
\let\subfigureautorefname=\figurename
\usepackage{adjustbox}
% Define block styles
\tikzset{%
  materia/.style={draw, fill=blue!20, text width=6.0em, text centered, minimum height=1.5em,drop shadow},
  etape/.style={materia, text width=8em, minimum width=10em, minimum height=3em, rounded corners, drop shadow},
  texto/.style={above, text width=6em, text centered},
  linepart/.style={draw, thick, color=black!50, -LaTeX, dashed},
  line/.style={draw, thick, color=black!50, -LaTeX},
  ur/.style={draw, text centered, minimum height=0.01em},
  back group/.style={fill=white!20,rounded corners, draw=black!50, dashed, inner xsep=15pt, inner ysep=10pt},
}
\pgfplotsset{%
  layers/standard/.define layer set={%
      background,axis background,axis grid,axis ticks,axis lines,axis tick labels,pre main,main,axis descriptions,axis foreground%
    }{
      grid style={/pgfplots/on layer=axis grid},%
      tick style={/pgfplots/on layer=axis ticks},%
      axis line style={/pgfplots/on layer=axis lines},%
      label style={/pgfplots/on layer=axis descriptions},%
      legend style={/pgfplots/on layer=axis descriptions},%
      title style={/pgfplots/on layer=axis descriptions},%
      colorbar style={/pgfplots/on layer=axis descriptions},%
      ticklabel style={/pgfplots/on layer=axis tick labels},%
      axis background@ style={/pgfplots/on layer=axis background},%
      3d box foreground style={/pgfplots/on layer=axis foreground},%
    },
}

\newcommand{\etape}[2]{node (p#1) [etape] {#2}}

\newcommand{\transreceptor}[3]{%
  \path [linepart] (#1.east) -- node [above] {\scriptsize #2} (#3);}
\tikzstyle{matheq} = [node distance=8.75cm, text width=21em, minimum width=1cm,
minimum height=2em, text centered,color=black!50]

\renewcommand{\baselinestretch}{1.2}

%
\newcommand{\gkc}{\|\gk\|^{1/2}}
\newcommand{\xk}{x_k}
\newcommand{\Hk}{H_k}
\newcommand{\gk}{g_k}
\newcommand{\dk}{d_k}
\newcommand{\gkn}{g_{k+1}}
\newcommand{\xkn}{x_{k+1}}
\newcommand{\dkn}{d_{k+1}}
\newcommand{\lk}{\lambda_{k+1}}
\newcommand{\cmark}{\ding{51}}%
\newcommand{\xmark}{\ding{55}}%
\newcommand{\arca}{\texttt{CubicReg-A}}

\newtheorem{theorem}{Theorem}[section]
\newtheorem{lemma}[theorem]{Lemma}

\theoremstyle{definition}
\newtheorem{definition}[theorem]{Definition}
\newtheorem{example}[theorem]{Example}
\newtheorem{xca}[theorem]{Exercise}

\theoremstyle{remark}
\newtheorem{remark}[theorem]{Remark}
\newtheorem{corollary}{Corollary}[section]
% \newtheorem{lemma}{Lemma}[section]
\newtheorem{proposition}{Proposition}[section]
\newtheorem{condition}{Condition}[section]
% \newtheorem{definition}{Definition}[section]
% \newtheorem{example}{Example}[section]
\newtheorem{question}{Question}[section]
% \newtheorem{remark}{Remark}[section]
\newtheorem{assumption}{Assumption}[section]
\newtheorem{strategy}{Strategy}[section]


\def\algorithmautorefname{Algorithm}
\def\subsectionautorefname{Section}
\def\definitionautorefname{Definition}
\def\corollaryautorefname{Corollary}
\def\lemmaautorefname{Lemma}
\def\assumptionautorefname{Assumption}
\def\propositionautorefname{Proposition}
\def\tableautorefname{Table}
\def\strategyautorefname{Strategy}
\def\remarkautorefname{Remark}
\def\conditionautorefname{Condition}


\newcommand{\hsodm}{\texttt{HSODM}}
\newcommand{\hsodmarc}{\texttt{HSODMArC}}
\newcommand{\hsodmhvp}{\texttt{HSODM-HVP}}
\newcommand{\drsom}{\texttt{DRSOM}}
\newcommand{\drsomh}{\texttt{DRSOM-H}}
\newcommand{\lbfgs}{\texttt{LBFGS}}
\newcommand{\newtontr}{\texttt{Newton-TR}}
\newcommand{\cg}{\texttt{CG}}
\newcommand{\arc}{\texttt{ARC}}
\newcommand{\atrms}{\texttt{NO-ATR}}
\newcommand{\itrace}{\texttt{TRACE}}
\newcommand{\gd}{\texttt{GD}}
\newcommand{\utr}{\texttt{UTR}}
\newcommand{\utrhvp}{\texttt{UTR-HVP}}
\newcommand{\newtontrst}{\texttt{Newton-TR-STCG}}
\newcommand{\iutr}{\texttt{iUTR(Hessian)}}
\newcommand{\iutrhvp}{\texttt{iUTR}}

\newcommand{\atr}{\texttt{GL-ATR}}

\makeatletter
\renewcommand\paragraph{\@startsection{paragraph}{4}%
  \z@ \z@ {-.5em}%
  {\normalfont\normalsize\bfseries}}%
\makeatother


\usepackage{float}
\usepackage{etoolbox}
\usepackage{cleveref}

% Define subroutine float
\newfloat{subroutine}{htbp}{los}
\floatname{subroutine}{Subroutine}

% Define cref name for subroutine
\crefname{subroutine}{subroutine}{subroutines}
\Crefname{subroutine}{Subroutine}{Subroutines}

\allowdisplaybreaks
% \usepackage[inline]{showlabels}
\title{Acceleration for Universal Trust-Region}
\author[1]{\small Yuntian Jiang\thanks{yuntianjiang07@163.sufe.edu.cn}}
\author[1]{\small Bo Jiang}
\author[1]{\small Chuwen Zhang}


\affil[1]{\footnotesize School of Information Management and Engineering\\ Shanghai University of Finance and Economics}

\begin{document}
\maketitle

% \section{Subproblem Properties}
% We now consider the accelerated version of the universal trust-region method proposed by \citet{jiang2023universal}. First, we go through the subproblem of UTR and establish some basic properties of it. In this paper, we consider a more general version:
% \begin{equation}\label{eq.newTR}
%     \begin{split}
%         \min_d ~&  \frac{1}{2}d^T\left (\nabla^2 f(\xk) + \sigma_k I \right ) d + \nabla f(\xk)^T d\\
%         \text{s.t. } ~& \|d\| \le r_k, \qquad \sigma_k, r_k > 0.
%     \end{split}
% \end{equation}
% we write the above subproblem for short as $\operatorname{UTR}(x_k,\sigma_k,r_k)$, the optimality condition of the above subproblem is 
% \begin{lemma}\label{lem.optimal condition}
%     The direction $\dk$ is the solution of \eqref{eq.newTR} if and only if there exists a dual multiplier $\lambda_k \geq 0 $ such that
%     \begin{subequations}
%         \begin{align}
%             \label{eq.optcond primal}
%              & \| d_k\| \leq r_k                                       \\
%             \label{eq.optcond coml slack}
%              & \lambda_k \left(\|d_k\|-r_k \right)=0                    \\
%             \label{eq.optcond firstorder}
%              & \left(\nabla^2 f(x_k) + \sigma_k I  + \lambda_k I \right) \dk = -\nabla f(x_k) \\
%             \label{eq.optcond secondorder}
%              & \nabla^2 f(x_k) + \sigma_k I + \lambda_k I \succeq 0.
%         \end{align}
%     \end{subequations}
% \end{lemma}
% From the above conditions, we have the following observation,
% \begin{lemma}
%     \label{eq.inner product g and d}
%     Suppose $d_k$ and $\lambda_k$ satisfy \eqref{eq.optcond primal}-\eqref{eq.optcond secondorder}, denote $\mu_k := \lambda_k+\sigma_k$, we have 
%     \begin{equation}
%         \label{eq.inner product g and d}
%         2\mu_k \langle \nabla f(x_k+d_k) ,-d_k\rangle \geq \|\nabla f(x_k+d_k)\|^2 +\mu_k^2\|d_k\|^2-\frac{M^2}{4}\|d_k\|^4.
%     \end{equation}
%     Further, if $\mu_k \geq \frac{1}{\sqrt{2}}M\|d_k\|$, we have 
%     \begin{equation}
%         \label{eq.inner product g and d further}
%         2\mu_k \langle \nabla f(x_k+d_k) ,-d_k\rangle \geq \|\nabla f(x_k+d_k)\|^2 +\frac{1}{2}\mu_k^2\|d_k\|^2.
%     \end{equation}
% \end{lemma}
% \begin{proof}
%     Due to the second-order Lipschitz continuity of $f$, we have 
%     \begin{equation*}
%         \|\nabla f(x_k+d_k)-\nabla f(x_k) -\nabla ^2 f(x_k) d_k\| \leq \frac{M}{2}\|d_k\|^2.
%     \end{equation*}
%     Plugging \eqref{eq.optcond firstorder} into the above equation, we have
%     \begin{equation*}
%         \|\nabla f(x_k+d_k)+\mu_k d_k\|\leq \frac{M}{2}\|d_k\|^2.
%     \end{equation*}
%     Squaring both sides we derive \eqref{eq.inner product g and d}.
% \end{proof}
% \section{Estimating Sequence}
% During the update of the algorithm, we intend to maintain the following two relations
% \begin{equation}
%     \label{eq.estimating sequence relation}
%     A_k f(x_k)\leq \phi_k^* :=\min \phi_k(x) \leq A_k f(x)+\frac{L}{2}\|x-x_0\|^2, \ \forall x\in \mathbb{R}^n.
% \end{equation}
% We define $\phi_0(x) = \frac{L}{2}\|x-x_0\|$, and $A_0=0$. Also, we recursively update $\phi_k$ and $A_k$ by 
% \begin{subequations}
%     \begin{align}
%     \label{eq.recursive A}
%         & A_{k+1} = A_k+a_{k+1}\\
%     \label{eq.recursive phi}
%         & \phi_{k+1}(x) = \phi_k(x)+a_{k+1} \left (f(z_k)+\langle \nabla f(z_k),x-z_k\rangle \right).
%     \end{align}
% \end{subequations}
% We now check the '$\leq$' relation, denote the minimizer of $\phi_k$ as $v_k$,
% \begin{equation*}
%     \begin{aligned}
%         \phi_{k+1}^* &= \min \left \{ \phi_k(x)+a_{k+1} \left (f(z_k)+\langle \nabla f(z_k),x-z_k\rangle \right)\right \}\\
%         &\geq \min \left \{ \phi_k^* +\frac{L}{2}\|x-v_k\|^2+a_{k+1} \left (f(z_k)+\langle \nabla f(z_k),x-z_k\rangle \right)\right \}\\
%         &\geq \min \left \{ A_k f(x_k) +\frac{L}{2}\|x-v_k\|^2+a_{k+1} \left (f(z_k)+\langle \nabla f(z_k),x-z_k\rangle \right)\right \}\\
%         &\geq \min \left \{ A_k f(z_k) +A_k \langle \nabla f(z_k),x_k-z_k\rangle +\frac{L}{2}\|x-v_k\|^2+a_{k+1} \left (f(z_k)+\langle \nabla f(z_k),x-z_k\rangle \right)\right \}\\
%         &= \min \left \{ A_{k+1} f(z_k)+A_{k+1} \langle \nabla f(z_k),\frac{A_k}{A_{k+1}}x_k+\frac{a_{k+1}}{A_{k+1}}v_k-z_k\rangle +\frac{L}{2}\|x-v_k\|^2+a_{k+1}\langle \nabla f(z_k),x-v_k\rangle\right \}\\
%         &= A_{k+1} f(z_k)+A_{k+1} \langle \nabla f(z_k),\frac{A_k}{A_{k+1}}x_k+\frac{a_{k+1}}{A_{k+1}}v_k-z_k\rangle-\frac{a_{k+1}^2}{2L}\|\nabla f(z_k)\|^2.
%     \end{aligned}
% \end{equation*}
% By letting $y_k = \frac{A_k}{A_{k+1}}x_k+\frac{a_{k+1}}{A_{k+1}}v_k$ and $x_{k+1} = \operatorname{UTR}(y_k,\sigma_k,r_k)$, we have 
% \begin{equation}
%     \label{eq.check relation 1}
%     \phi_{k+1}^* \geq A_{k+1} f(x_{k+1})+ A_{k+1}\langle \nabla f(x_{k+1}),y_k-x_{k+1}\rangle-\frac{a_{k+1}^2}{2L}\|\nabla f(x_{k+1})\|^2.
% \end{equation}
% The difficulty lies in that we have to guarantee 
% $$A_{k+1}\langle \nabla f(x_{k+1}),y_k-x_{k+1}\rangle-\frac{a_{k+1}^2}{2L}\|\nabla f(x_{k+1})\|^2\geq 0$$ 
% while keeping $a_{k+1}$ as large as possible. These quantities seem to link with the multiplier $\mu_k$, which is posterior. When $\mu_k > \sigma_k$, i.e., $\lambda_k>0$, relation 2 may break down, we can apply a damping technique, define 
% \begin{gather}
%     \label{eq.damping parameter}
%     \gamma_{k+1}:= \frac{\sigma_k}{\mu_k}, \ a'_{k+1}:=\gamma_{k+1}a_{k+1}, \ A'_{k+1} = A_k+a'_{k+1},\\
%     \label{eq.damping x}
%     x'_{k+1} =\frac{(1-\gamma_{k+1})A_k}{A'_{k+1}}x_k+\frac{\gamma_{k+1}A_{k+1}}{A'_{k+1}}x_{k+1}. 
%     \end{gather}
%     Then we can maintain the relation
%     \begin{equation}
%         \label{eq.damped relation}
%         A'_{k+1}f(x'_{k+1}) \leq \phi'^*_{k+1} =\min \phi'_{k+1}(x) := \phi_k(x)+a'_{k+1} \left ( f(x_{k+1}+\langle \nabla f(x_{k+1}),x-x_{k+1}\rangle\right ).
%     \end{equation}
%     Or we can try another kind of damping, we set 
%     \begin{equation*}
%         \phi_{k+1} := \gamma_k \phi_k +a_k\left(f(z_k)+\langle \nabla f(z_k),x-z_k\rangle \right), \ A_{k+1} = \gamma_k A_k +a_{k}.
%     \end{equation*}
%     Then we have the following recursive relation:
%     \begin{equation*}
%         \begin{aligned}
%             \phi_{k+1}(x) &= \gamma_k \phi_k(x)+a_k f(z_k)+a_k \langle \nabla f(z_k),x-z_k\rangle\\
%             &\geq \gamma_k \phi_k^* +\gamma_k \sigma_k \|x-v_k\|^2+ a_k f(z_k)+a_k \langle \nabla f(z_k),x-z_k\rangle\\
%             &\geq \gamma_k A_k f(x_k)+\gamma_k \sigma_k \|x-v_k\|^2+ a_k f(z_k)+a_k \langle \nabla f(z_k),x-z_k\rangle\\
%             &= A_{k+1} f(z_k)+A_{k+1}\langle \nabla f(z_k),y_k-z_k\rangle+\gamma_k\sigma_k \|x-v_k\|^2+ a_k \langle \nabla f(z_k),x-v_k\rangle.
%         \end{aligned}
%     \end{equation*}
%     Set $z_k = x_{k+1}$ and minimize both sides with respect to $x$, we have 
%     \begin{equation*}
%         \phi_{k+1}^* \geq A_{k+1} f(x_{k+1}) +A_{k+1}\langle \nabla f(x_{k+1}),y_k-x_{k+1} \rangle -\frac{a_k^2}{4\gamma_k \sigma_k}\|\nabla f(x_{k+1})\|^2.
%     \end{equation*}
%     we give an informal algorithm as below:
%     \begin{algorithm}[ht]
%     \caption{An Accelerated Trust-Region Method}\label{alg.accelerated utr}
%     \begin{algorithmic}[1]
%         \STATE{\textbf{input:} $x_0,v_0\in \mathbb{R}^n$, $A_0=0$, $\gamma_0=0$;}
%         \FOR{$k=0, 1, \ldots, \infty$}
%         \STATE{$a_k = \frac{1+\sqrt{1+4\gamma_k A_k}}{2}$}
%         \STATE{$y_k = \frac{A_k\gamma_k}{A_k\gamma_k+a_k}x_k+\frac{a_k}{A_k\gamma_k+a_k}v_k$}
%         \STATE{$x_{k+1} = \operatorname{UTR}(y_k,\frac{\sqrt{M}}{3}\|\nabla f(x_k)\|^{1/2},\frac{1}{3\sqrt{M}}\|\nabla f(x_k)\|^{1/2})$;}
%         \STATE{$A_{k+1} = A_k\gamma_k+a_k$;}
%         \STATE{$v_{k+1} = v_k-\frac{a_{k+1}}{\gamma_k}\nabla f(x_{k+1})$;}
%         \STATE{$\gamma_{k+1} = \frac{\|\nabla f(x_{k+1})\|^{1/2}}{\|\nabla f(x_k)\|^{1/2}}$;}
%         \ENDFOR
%     \end{algorithmic}
% \end{algorithm}

% Or we can directly set $\phi_{k+1}(x):= A_{k+1}f(x_{k+1})+\frac{\sigma_{k+1}}{2}\|x-v_{k+1}\|^2$, where $a_{k+1}$, $v_{k+1}$ can be set with freedom. Then the relation $\phi_{k+1}^* \geq A_{k+1}f(x_{k+1})$ is automatically guaranteed, then we needs to check 
% \begin{equation*}
%     \phi_k(x) \leq A_{k+1}f(x)+ h_{k+1}(x),
% \end{equation*}
% ideally, $h$ is also a quadratic or cubic function.
% \clearpage
\section{Introduction}
In this paper, we consider the following unconstrained optimization problem:
\begin{equation}
    \label{eq.main prob}
    \min_{x\in\mathbb{R}^n} \ f(x),
\end{equation}
where $f$ is convex, twice continuously differentiable, and bounded below by $f^*$. We aim to find the approximate solutions of \eqref{eq.main prob}, the approximate solutions should satisfy the following
\begin{definition}
    \label{def.convex approximate solution}
    For $0<\epsilon<1$, $x$ is called an $\epsilon$-approximate solution of $f$ if 
    \begin{equation}
        \label{eq.convex approximate solution}
        f(x)-f^* \leq \epsilon \ \text{or} \ \|\nabla f(x)\| \leq \epsilon.
    \end{equation}
\end{definition}
\section{Preliminaries}
We now consider the accelerated version of the universal trust-region method proposed by \citet{jiang2023universal}. We still adopt the commonly used Hessian Lipschitz and bounded assumption:
\begin{assumption}
    \label{assm.bounded}
    The Hessian $\nabla^2 f(x)$ of the objective function is bounded by constant $\kappa_H>0$, i.e.,
    \begin{equation}
        \label{eq.bounded hessian}
        \|\nabla^2 f(x)\| \leq \kappa_H \quad \forall x\in \mathbb{R}^n.
    \end{equation}
\end{assumption}

\begin{assumption}
    \label{assm.lipschitz}
    The Hessian $\nabla^2 f(x)$ of the objective function is Lipschitz continuous with constant $M>0$, i.e.,
    \begin{equation}
        \label{eq.lipschiz}
        \|\nabla^2 f(x)-\nabla^2 f(y)\| \leq M\|x-y\| \quad \forall x,y\in\mathbb{R}^n.
    \end{equation}
\end{assumption}
As a direct consequence, we have the following lemma:
\begin{lemma}[Lemma 4.1.1, \citet{nesterov_lectures_2018}]\label{lem.lipschitz}
    If \(f:\mathbb{R}^n \mapsto \mathbb{R}\) satisfies \autoref{assm.lipschitz}, then for all \(x,y\in \mathbb{R}^n\), we have
    \begin{subequations}
        \begin{align}
            \label{eq.first-order exp}
             & \left\|\nabla f(y)-\nabla f(x)-\nabla^{2} f(x)(y-x)\right\|                       \leq \frac{M}{2}\|y-x\|^{2} \\
            \label{eq.second-order exp}
             & \left|f(y)-f(x)-\nabla f(x)^T(y-x)-\frac{1}{2}(y-x)^T\nabla^{2} f(x)(y-x)\right|  \leq \frac{M}{6}\|y-x\|^{3}
        \end{align}
    \end{subequations}
\end{lemma}

First, we go through the subproblem of UTR and establish some basic properties of it. In this paper, we consider a more general version, at point $x$, the UTR subproblem with parameter $\sigma,r>0$ is
\begin{equation}\label{eq.newTR}
    \begin{split}
        \min_d ~&  \frac{1}{2}d^T\left (\nabla^2 f(x) + \sigma I \right ) d + \nabla f(x)^T d\\
        \text{s.t. } ~& \|d\| \le r, \qquad \sigma, r > 0.
    \end{split}
\end{equation}
The optimality condition of the above subproblem is 
\begin{lemma}\label{lem.optimal condition}
    The direction $d$ is the solution of \eqref{eq.newTR} if and only if there exists a dual multiplier $\lambda \geq 0 $ such that
    \begin{subequations}
        \begin{align}
            \label{eq.optcond primal}
             & \| d\| \leq r                                       \\
            \label{eq.optcond coml slack}
             & \lambda \left(\|d\|-r \right)=0                    \\
            \label{eq.optcond firstorder}
             & \left(\nabla^2 f(x) + \sigma I  + \lambda I \right) d = -\nabla f(x) \\
            \label{eq.optcond secondorder}
             & \nabla^2 f(x) + \sigma I + \lambda I \succeq 0.
        \end{align}
    \end{subequations}
\end{lemma}
We write the exact primal-dual solution pair of \eqref{eq.newTR} for short as $d,\lambda=\operatorname{UTR}(x,\sigma,r)$, we note here that the $\operatorname{UTR}$ oracle is the main oracle that we use, we will evaluate the complexity of an algorithm in terms of the number of the oracles called to find the approximate solutions defined in \eqref{eq.convex approximate solution}. From the above conditions, we have the following observations, the first one is with respect to the decrease of the objective function value:
\begin{lemma}
    \label{lem.function decrease}
    Suppose \autoref{assm.lipschitz} holds and $d_k,\lambda_k = \operatorname{UTR}(x_k,\sigma_k,r_k)$, then we have
    \begin{equation}
        \label{eq.function decrease}
        f(x_k+d_k) \leq f(x_k) -\frac{1}{2}d_k^T \nabla^2 f(x_k) d_k -(\lambda_k+\sigma_k)\|d_k\|^2 +\frac{M}{6}\|d_k\|^3.
    \end{equation}
    Especially, when $\sigma_k \geq M r_k$ and $\lambda_k>0$, we have
    \begin{equation}
        \label{eq.function decrease advanced}
        f(x_k+d_k) \leq f(x_k)-\frac{1}{2}d_k^T \nabla^2 f(x_k) d_k-\lambda_k r_k^2 -\frac{5M}{6} r_k^3.
    \end{equation}
\end{lemma}
\begin{proof} 
    Set $y=x_k+d_k$, $x=x_k$ in \eqref{eq.second-order exp}, we have 
    \begin{equation*}
        f(x_k+d_k) \leq f(x_k) +\nabla f(x_k)^T d_k +\frac{1}{2}d_k^T \nabla ^2 f(x_k) d_k +\frac{M}{6}\|d_k\|^3.
    \end{equation*}
    Plug \eqref{eq.optcond firstorder} into the above we have \eqref{eq.function decrease}. Note that when $\lambda_k>0$, from \eqref{eq.optcond coml slack} we have $\|d_k\| = r_k$, from \eqref{eq.function decrease}, we have
    \begin{equation*}
        \begin{aligned}
            f(x_k+d_k)&\leq f(x_k)-\frac{1}{2}d_k^T \nabla^2 f(x_k)d_k-\lambda_k r_k^2-\sigma_k r_k^2+\frac{M}{6}r_k^3\\
            &\leq f(x_k)-\frac{1}{2}d_k^T \nabla^2 f(x_k) d_k-\lambda_k r_k^2 -\frac{5M}{6} r_k^3.
        \end{aligned}
    \end{equation*}
    In the second line, we used $\sigma_k \geq M r_k.$
\end{proof}
\begin{corollary}
    Suppose \autoref{assm.lipschitz} holds, $d_k,\lambda_k=\operatorname{UTR}(x_k,\sigma_k=\frac{\sqrt{M}}{3} \|\nabla f(x_k)\|^{1/2},r_k= \frac{1}{3\sqrt{M}} \|\nabla f(x_k)\|^{1/2})$ and $\lambda_k>0$, we have
    \begin{equation}
        \label{eq.decrease of order 1.5}
        f(x_k+d_k) \leq f(x_k)-\frac{5}{162\sqrt{M}}\|\nabla f(x_k)\|^{3/2}.
    \end{equation}
\end{corollary}
As for the bound for the gradient, we have
\begin{lemma}
    \label{lem.bound next gradient}
    Suppose \autoref{assm.lipschitz} holds and $d_k,\lambda_k= \operatorname{UTR}(y_k,\sigma_k,r_k)$ , then we have
    \begin{equation}
        \label{eq.bound next gradient}
        \|\nabla f(x_k+d_k)\| \leq \frac{M}{2}r_k^2+(\sigma_k+\lambda_k)r_k.
    \end{equation}
\end{lemma}
\begin{proof}
    From \eqref{eq.second-order exp}, we have
    \begin{equation*}
        \|\nabla f(x_k+d_k)- \nabla f(x_k) -\nabla ^2 f(x_k) d_k\| \leq \frac{M}{2}\|d_k\|^2.
    \end{equation*}
    By \eqref{eq.optcond firstorder} and the above, we have
    \begin{equation*}
        \|\nabla f(x_k+d_k)+(\sigma_k+\lambda_k) d_k\| \leq \frac{M}{2}\|d_k\|^2.
    \end{equation*}
    Note $\|d_k\|\leq r_k$, apply the triangle inequality and rearrange items we can derive \eqref{eq.bound next gradient}.
\end{proof}
The following lemma is a crucial lemma for a second-order algorithm to be accelerated.
\begin{proof}
    Here $\sigma_k \geq M r_k$ holds, hence we can directly apply \autoref{lem.function decrease} with $r_k = \frac{1}{3\sqrt{M}} \|\nabla f(x_k)\|^{1/2}.$
\end{proof}
\begin{lemma}
    \label{lem.inner product g and d}
    Suppose \autoref{assm.lipschitz} holds, and 
    \begin{equation}
        \label{eq.regularized condition for d}
        \left (\nabla^2 f(x_k)+\mu_k I\right)d_k = -\nabla f(x_k),
    \end{equation}
    holds for some $\mu_k$ with $\mu_k\geq M \|d_k\|$, then we have
    \begin{equation}
        \label{eq.inner product g d}
        \langle \nabla f(x_k+d_k),-d_k \rangle \geq \frac{\|\nabla f(x_k+d_k)\|^2}{2\mu_k}+\frac{3}{8}\mu_k \|d_k\|^2.
    \end{equation}
    Further, if $\mu_k \leq 2M\|d_k\|$, we have
    \begin{equation}
        \label{eq.inner product g d advanced}
        \langle \nabla f(x_k+d_k),-d_k \rangle \geq \frac{\sqrt{6}}{6} \frac{\|\nabla f(x_k+d_k)\|^{3/2}}{\sqrt{M}}.
    \end{equation}
\end{lemma}
\begin{proof}
    Set $y=x_k+d_k$, $x= x_k$ in \eqref{eq.first-order exp}, we have 
    \begin{equation*}
        \|\nabla f(x_k+d_k)-\nabla f(x_k)-\nabla ^2 f(x_k) d_k\| \leq \frac{M}{2}\|d_k\|^2.
    \end{equation*}
    Plug \eqref{eq.regularized condition for d} into the above, we have
    \begin{equation*}
        \|\nabla f(x_k+d_k)+\mu_k d_k\| \leq \frac{M}{2}\|d_k\|^2.
    \end{equation*}
    Squaring both sides and rearranging items, we have
    \begin{equation*}
    \begin{aligned}
        2\mu_k \langle \nabla f(x_k+d_k),-d_k \rangle &\geq \|\nabla f(x_k+d_k)\|^2 +\mu_k^2 \|d_k\|^2 -\frac{M}{4}\|d_k\|^4\\
        &\geq \|\nabla f(x_k+d_k)\|^2 +\frac{3}{4}\mu_k^2 \|d_k\|^2.
    \end{aligned}
    \end{equation*}
    The second line is due to $\mu_k\geq M\|d_k\|$, dividing both sides by $2\mu_k$ gives \eqref{eq.inner product g d}. Further, when $\mu_k \leq 2M\|d_k\|$, the above gives
    \begin{equation*}
        4M\|d_k\| \langle \nabla f(x_k+d_k),-d_k \rangle \geq \|f(x_k+d_k)\|^2 +\frac{3}{4} M^2 \|d_k\|^4,
    \end{equation*}
    which yields
    \begin{equation*}
        \langle \nabla f(x_k+d_k),-d_k \rangle \geq \frac{\|f(x_k+d_k)\|^2}{4M\|d_k\|}+\frac{3}{16}M\|d_k\|^3.
    \end{equation*}
    Consider an auxiliary function $h(t) = \frac{\|f(x_k+d_k)\|^2}{4Mt}+\frac{3}{16}Mt^3$ where $t\geq 0$, taking derivatives gives $h(t)$ achieves its minimum at $t^*= \frac{\sqrt{2}\|\nabla f(x_k+d_k)\|^{1/2}}{\sqrt{3M}}$, plugging $t^*$ back gives \eqref{eq.inner product g d advanced}.
\end{proof}

\section{The Nesterov Acceleration}
    We give the first accelerated \autoref{alg.accelerated utr} and its auxiliary search procedure \autoref{alg.bisection} as below:
    \begin{algorithm}[ht]
    \caption{Accelerated Trust-Region(ATR) Method}\label{alg.accelerated utr}
    \begin{algorithmic}[1]
        \STATE{\textbf{input:} $x_0=v_0\in \mathbb{R}^n$, $s_0=0$;}
        \FOR{$k=0, 1, \ldots, \infty$}
        \STATE{$y_k = \frac{k}{k+3}x_k+\frac{3}{k+3}v_k$}
        \STATE{$(d_k,\lambda_k) = \operatorname{UTR}(y_k,\frac{\sqrt{M}}{3}\|\nabla f(y_k)\|^{1/2},\frac{1}{3\sqrt{M}}\|\nabla f(y_k)\|^{1/2})$}
        \IF{$\lambda_k>0$}
        \STATE{$x_{k+1}=y_k+d_k$}
        \STATE{$s_k = s_{k-1}+\frac{(k+1)(k+2)}{2}\nabla f(y_k)$}
        \ELSE
        \STATE{$d_k^{(0)}=d_k$, $\lambda_k^{(0)} = \lambda_k$, $j=0$}
        \WHILE{$\lambda_k^{(j)}=0$}
        \STATE{$j=j+1$}
        \STATE{$(d_k^{(j)},\lambda_k^{(j)})=\operatorname{UTR}(y_k,\frac{\sqrt{M}}{3}\|\nabla f(y_k+d_k^{(j-1)})\|^{1/2},\frac{1}{3\sqrt{M}}\|\nabla f(y_k+d_k^{(j-1)})\|^{1/2})$}
        \ENDWHILE
        \IF{$\lambda_k^{(j)}\leq \frac{\sqrt{M}}{3}\|\nabla f(y_k+d_k^{(j-1)})\|^{1/2}$}
        % \STATE{$x_{k+1} = y_k+d_k^{(j)}$, $\mu_k = \lambda_k^{(j)}+\frac{\sqrt{M}}{3}\|\nabla f(y_k+d_k^{(j-1)})\|^{1/2}$}
        \STATE{$x_{k+1} = y_k+d_k^{(j)}$}
        \label{line.adaptive search output}
        \ELSE
        \STATE{$\mu_- = \frac{\sqrt{M}}{3}\|\nabla f(y_k+d_k^{(j-1)})\|^{1/2}$, $\mu_+ = \lambda_k^{(j)}+\frac{\sqrt{M}}{3}\|\nabla f(y_k+d_k^{(j-1)})\|^{1/2}$}
        % \STATE{$x_{k+1},\mu_k = \operatorname{BISECTION\_ATR}(y_k,\mu_-,\mu_+)$}
        \STATE{$x_{k+1} = \operatorname{BISECTION\_ATR}(y_k,\mu_-,\mu_+)$}
        \ENDIF
        \STATE{$s_k=s_{k-1}+\frac{(k+1)(k+2)}{2}\nabla f(x_{k+1})$}
        \ENDIF
        \STATE{$v_{k+1} = v_0-\sqrt{\frac{\ell}{3M\|s_k\|}}s_k$}
        \ENDFOR
    \end{algorithmic}
\end{algorithm}

\begin{algorithm}[!ht]
    \caption{Bisection for ATR}\label{alg.bisection}
    \begin{algorithmic}[1]
        \STATE{\textbf{input:} $y\in \mathbb{R}^n$, $\mu_- < \mu_+$;}
        \WHILE{True}
        \STATE{$\mu = \frac{\mu_-+\mu_+}{2}$}
        \STATE{$d = -\left (\nabla ^2 f(y)+\mu I\right )^{-1}\nabla f(y)$}
        \IF{$\frac{\mu}{\|d\|}<M$}
        \STATE{$\mu_-=\mu$}
        \ELSIF{$\frac{\mu}{\|d\|}>2M$}
        \STATE{$\mu_+ = \mu$}
        \ELSE
        \STATE{\textbf{output} $y+d,\mu$}
        \ENDIF
        \ENDWHILE
    \end{algorithmic}
\end{algorithm}
We first make some clarifications to show that \autoref{alg.accelerated utr} is well-defined. In the $k$-th outer loop of \autoref{alg.accelerated utr}, there is an inner loop indexed by $j$, we first show that the inner loop will terminate eventually,
\begin{lemma}
    \label{lem.inner loop bound next gradient}
    In $k$-th outer loop and $j$-th inner loop, if $\lambda_k^{(j)}=0$, then we have
    \begin{equation}
        \label{eq.inner loop bound next gradient}
        \|\nabla f(y_k+d_k^{(j)}) \| \leq \frac{1}{18}\|\nabla f(y_k+d_k^{(j-1)})\|.
    \end{equation}
\end{lemma}
\begin{proof}
    This lemma is implied by \autoref{lem.bound next gradient} with $\sigma_k = \frac{\sqrt{M}}{3}\|\nabla f(y_k+d_k^{(j-1)})\|$, $\lambda_k = 0$.
\end{proof}
Denote $j_k$ as the number of the inner loop called in the $k$-th outer loop, for those index $k$ with $\lambda_k>0$, the inner loop is not called, we denote $j_k=0$. Therefore, we can classify the outer loop index $k$ into two categories:
\begin{equation}
    \label{eq.outer index catagory}
    \mathcal{Z} = \left \{ k | j_k=0\right\}, \mathcal{N} = \left \{k | j_k >0\right\}.
\end{equation}
Also, for ease of analysis, we define the following two sets with respect to outer index $p$:
\begin{equation}
    \label{eq.outer index before p}
    \mathcal{Z}_{p} = \left \{ k<p | j_k=0\right\}, \mathcal{N}_{p} = \left \{k<p | j_k >0\right\}.
\end{equation}
We now show that $j_k$ is finite:
\begin{corollary}
    \label{coro.inner loop terminate}
    In the $k$-th outer loop of \autoref{alg.accelerated utr}, the inner loop will terminate in 
    \begin{equation*}
        j_k \leq O\left (\log \frac{\|\nabla f(y_k)\|}{\epsilon}\right)
    \end{equation*}
    iterations, otherwise we find a point such that
    \begin{equation*}
        \| \nabla f(y_k+d_k^{(j_k)})\| \leq \epsilon.
    \end{equation*}
    Since the inner loop will terminate eventually, we can actually have a tighter bound for $j_k$:
    \begin{equation}
        \label{eq.tighter bound jk}
        j_k \leq O\left (\log \frac{\|\nabla f(y_k)\|}{\|\nabla f(y_k+d_k^{(j_k-1)})\|}\right)
    \end{equation}
\end{corollary}

\subsection{Global Convergence of \autoref{alg.accelerated utr}}
In this section, we investigate the complexity of \autoref{alg.accelerated utr} for finding an approximate solution as in \autoref{def.convex approximate solution}. We analyze the convergence of \autoref{alg.accelerated utr} by the estimating sequence technique \cite{nesterov_lectures_2018}. We first recursively define the following auxiliary sequences for $k\in \mathbb{N}$,
\begin{equation}
    \label{eq.update a and A}
    a_k = \frac{(k+1)(k+2)}{2}, \ A_k=\frac{k(k+1)(k+2)}{6}
\end{equation}
\begin{align}\label{eq.update phi}
\phi_0(x) = \frac{M}{\ell}\|x-x_0\|^3, \ \phi_{k+1}(x) =\left\{ 
\begin{aligned}
&\phi_k(x)+a_k\left (f(x_k)+\langle \nabla f(x_k),x-x_k\rangle \right ), k\in \mathcal{N}\\
&\phi_k(x)+a_k\left (f(y_k)+\langle \nabla f(y_k),x-y_k\rangle \right ), k\in \mathcal{Z}
\end{aligned}
\right.
\end{align}

Here are some basic properties of the estimating sequence, which is modified from \citet{nesterov_lectures_2018}.
\begin{lemma}
    \label{lem.magnitude a and A}
    For $k\geq 0$, we have
    \begin{equation}
        \label{eq.magnitude a and A}
        \frac{A_{k+1}}{a_k^{3/2}} \geq \frac{\sqrt{2}}{3}.
    \end{equation}
\end{lemma}
\begin{proof}
    $\frac{A_{k+1}}{a_k^{3/2}}= \frac{(k+1)(k+2)(k+3)}{6}\cdot \frac{2^{3/2}}{(k+1)^{3/2}(k+2)^{3/2}} = \frac{\sqrt{2}(k+3)}{3(k+1)^{1/2}(k+2)^{1/2}} \geq \frac{\sqrt{2}}{3}.$
\end{proof}
We call a differentiable function $d(x)$ on $\mathbb{R}^n$ uniformly convex \citet[Section 4.2.2]{nesterov_lectures_2018} of degree $p\geq 2$ with constant $q>0$ if
\begin{equation}
    \label{eq.uniformly convex of p}
    d(y)\geq d(x)+\langle \nabla d(x),y-x\rangle +\frac{q}{p}\|y-x\|^p, \ \forall x,y\in\mathbb{R}^n.
\end{equation}
\begin{lemma}
    \label{lem.property of phi}
    For the function sequence $\phi_k(x)$, they have the following properties
    \begin{enumerate}
        \item $\phi_k(x)$ is uniformly convex of degree $3$ with constant $\frac{3M}{2\ell}$.
        \item $v_k$ is the unique minimizer of $\phi_k(x)$, and 
        \begin{equation}
            \label{eq.phi_k phi_low}
            \phi_k(x)\geq \phi_k^* + \frac{M}{2\ell}\|x-v_k\|^3.,
        \end{equation}
        where $\phi_k^* =\min \phi_k(x).$
    \end{enumerate}
\end{lemma}
\begin{proof}
    The first claim is from the definition of $\phi_0$ in \eqref{eq.update phi}, for analysis of uniform convex functions, please refer to \citet[Section 4.2.2]{nesterov_lectures_2018}. For the second claim, we have
    \begin{equation}
        \label{eq.phi recursive}
        \phi_k(x) = \phi_0(x)+\sum_{i\in \mathcal{N}_{k}}a_i\left (f(y_i)+\langle \nabla f(y_i),x-y_i\rangle \right)+\sum_{j\in \mathcal{Z}_{k}}a_j\left (f(x_j)+\langle \nabla f(x_j),x-x_j\rangle \right).
    \end{equation}
    From optimality condition
    \begin{equation*}
    \begin{aligned}
        \nabla \phi_k(x) &= \nabla \phi_0(x)+\sum_{i\in \mathcal{N}_k}a_i\nabla f(y_i)+ \sum_{j\in \mathcal{Z}_k}a_j\nabla f(x_j)\\
        &= \frac{3M}{\ell}\|x-x_0\|(x-x_0)+\sum_{i\in \mathcal{N}_k}a_i\nabla f(y_i)+ \sum_{j\in \mathcal{Z}_k}a_j\nabla f(x_j)\\
        &=0.
    \end{aligned}
    \end{equation*}
    Solving the optimality condition and noticing the way we update $s_k$ gives 
    \begin{equation*}
        x = v_0 - \sqrt{\frac{\ell}{3M\|s_{k-1}\|}}s_{k-1},
    \end{equation*}
    which is $v_k$. Then \eqref{eq.phi_k phi_low} comes from the factor that $\phi_k(x)$ is uniformly convex of degree 3.
\end{proof}
During the update of the algorithm, we can show that the following properties hold
\begin{lemma}\label{lem.estimate sequence}
If we set $0<\ell \leq\frac{25}{12\cdot 54^2}$, then we can guarantee the following for all $k\geq 0$:
\begin{equation}
    \label{eq.estimating sequence relation}
    A_k f(x_k)\leq \phi_k^*  \leq \phi_k(x)\leq A_k f(x)+\phi_0(x), \ \forall x\in \mathbb{R}^n.
\end{equation}
\end{lemma}
\begin{proof}
We prove by induction, note that $A_0 = 0$, so \eqref{eq.estimating sequence relation} holds for $i=0$. Suppose it also holds for $i=k$, 
We first check  
$$\phi_{k+1}(x)\leq A_{k+1} f(x)+\phi_0(x),$$ 
suppose $k\in \mathcal{N}$
\begin{equation*}
    \begin{aligned}
        \phi_{k+1}(x)&=\phi_k(x)+a_k\left (f(y_k)+\langle \nabla f(y_k),x-y_k\rangle \right)\\
        &\leq A_k f(x)+\phi_0(x) +a_k\left (f(y_k)+\langle \nabla f(y_k),x-y_k\rangle \right)\\
        & \leq A_k f(x)+\phi_0(x) +a_k f(x)\\
        &=A_{k+1} f(x)+\phi_0(x).
    \end{aligned}
\end{equation*}
The first inequality is from the induction hypothesis. The proof for the case $k\in \mathcal{Z}$ is exactly the same.

We now check $\phi_{k+1}^* \geq A_{k+1}f(x_{k+1})$,
\begin{equation*}
    \begin{aligned}
        \phi_{k+1}^* &= \min \left \{ \phi_k(x)+a_k \left (f(z_k)+\langle \nabla f(z_k),x-z_k\rangle \right)\right \}\\
        &\geq \min \left \{ \phi_k^* +\frac{M}{2\ell}\|x-v_k\|^3+a_k \left (f(z_k)+\langle \nabla f(z_k),x-z_k\rangle \right)\right \}\\
        &\geq \min \left \{ A_k f(x_k) +\frac{M}{2\ell}\|x-v_k\|^3+a_k \left (f(z_k)+\langle \nabla f(z_k),x-z_k\rangle \right)\right \}\\
        &\geq \min \left \{ A_k f(z_k) +A_k \langle \nabla f(z_k),x_k-z_k\rangle +\frac{M}{2\ell}\|x-v_k\|^3+a_k \left (f(z_k)+\langle \nabla f(z_k),x-z_k\rangle \right)\right \}\\
        &= \min \left \{ A_{k+1} f(z_k)+A_{k+1} \langle \nabla f(z_k),\frac{A_k}{A_{k+1}}x_k+\frac{a_k}{A_{k+1}}v_k-z_k\rangle +\frac{M}{2\ell}\|x-v_k\|^3+a_k\langle \nabla f(z_k),x-v_k\rangle\right \}\\
        &= A_{k+1} f(z_k)+A_{k+1} \langle \nabla f(z_k),\frac{A_k}{A_{k+1}}x_k+\frac{a_k}{A_{k+1}}v_k-z_k\rangle-\frac{2\sqrt{2}\sqrt{\ell}\cdot a_k^{3/2}}{3\sqrt{3}\cdot\sqrt{M}}\|\nabla f(z_k)\|^{3/2}.
    \end{aligned}
\end{equation*}
In the above analysis, the first line is because of the uniform convexity of $\phi_k$, the second line is from the induction hypothesis, and the third line is from the convexity of $f$. Note that we set $z_k = y_k$ if $k\in \mathcal{N}$, otherwise, $z_k=x_k$. When $k\in \mathcal{N}$, we have
\begin{equation*}
\begin{aligned}
    \phi_{k+1}^* &\geq A_{k+1}f(y_k)-\frac{2\sqrt{2}\sqrt{\ell}\cdot a_k^{3/2}}{3\sqrt{3}\cdot\sqrt{M}}\|\nabla f(y_k)\|^{3/2}\\
    &\geq f(x_{k+1}).
\end{aligned}
\end{equation*}
The second inequality is from the way we set $\ell$, \eqref{eq.decrease of order 1.5}, and \eqref{eq.magnitude a and A}. 

When $k\in \mathcal{Z}$, we have 
\begin{equation*}
\begin{aligned}
    \phi_{k+1}^* &\geq A_{k+1} f(x_{k+1}) +A_{k+1}\langle \nabla f(x_{k+1}),y_k-x_{k+1} \rangle -\frac{2\sqrt{2}\sqrt{\ell}\cdot a_k^{3/2}}{3\sqrt{3}\cdot\sqrt{M}}\|\nabla f(x_{k+1})\|^{3/2}\\
    &= A_{k+1} f(x_{k+1}) +A_{k+1}\left (\langle \nabla f(x_{k+1}),y_k-x_{k+1} \rangle -A_{k+1}^{-1}\frac{2\sqrt{2}\sqrt{\ell}\cdot a_k^{3/2}}{3\sqrt{3}\cdot\sqrt{M}}\|\nabla f(x_{k+1})\|^{3/2}\right )\\
    &\geq A_{k+1} f(x_{k+1}) +A_{k+1}\left (\langle \nabla f(x_{k+1}),y_k-x_{k+1} \rangle -\frac{2\sqrt{\ell}}{\sqrt{3}\cdot\sqrt{M}}\|\nabla f(x_{k+1})\|^{3/2}\right ).
    \end{aligned}
\end{equation*}
The last inequality is from \eqref{eq.magnitude a and A}. Now it reduces to prove 
\begin{equation*}
    \langle f(x_{k+1}),y_k-x_{k+1} \rangle -\frac{2\sqrt{\ell}}{\sqrt{3}\cdot\sqrt{M}}\|\nabla f(x_{k+1})\|^{3/2} \geq 0.
\end{equation*}
Note that the output of line \autoref{line.adaptive search output} of \autoref{alg.accelerated utr} and the bisection procedure satisfies
\begin{gather*}
    x_{k+1}=y_k+d_k\\
    \left (\nabla ^2 f(y_k)+\mu_k I\right)d_k = -\nabla f(y_k)\\
    M\|d_k\| \leq \mu_k \leq 2M\|d_k\|.
\end{gather*}
As a result of \eqref{eq.inner product g d advanced}, we have 
\begin{equation*}
    \langle f(x_{k+1}),y_k-x_{k+1} \rangle -\frac{2\sqrt{\ell}}{\sqrt{3}\cdot\sqrt{M}}\|\nabla f(x_{k+1})\|^{3/2} \geq \frac{\sqrt{6}}{6\sqrt{M}}\|\nabla f(x_{k+1})\|^{3/2}-\frac{2\sqrt{\ell}}{\sqrt{3}\cdot\sqrt{M}}\|\nabla f(x_{k+1})\|^{3/2} \geq 0.
\end{equation*}
Thus the proof is finished.
\end{proof}
As a direct consequence, we have the following convergence result of \autoref{alg.accelerated utr},
\begin{theorem}
    \label{thm.convergence outer loop}
    Suppose \autoref{assm.bounded} and \autoref{assm.lipschitz} hold, then \autoref{alg.accelerated utr} find a point $x\in\mathbb{R}^n$ satisfies \eqref{eq.convex approximate solution} in $O(\epsilon^{-1/3})$ outer iterations.
\end{theorem}
\begin{proof}
    This theorem is easily derived by combining \autoref{lem.estimate sequence} and \autoref{coro.inner loop terminate}.
\end{proof}
\subsection{Bisection Procedure}
Note that \autoref{alg.accelerated utr} relies on a bisection procedure, we now conduct the complexity analysis for this procedure.


Suppose \autoref{alg.accelerated utr} enters the bisection procedure at $k$-th outer iteration and $j_k$-th inner iteration, then by the mechanism of \autoref{alg.accelerated utr}, the following can be established:
    \begin{gather}
    \label{eq.magnitude lambda}
        \lambda_k^{(j_k)} > \frac{\sqrt{M}}{3}\|\nabla f(y_k+d_k^{(j_k-1)})\|^{1/2}\\
    \label{eq.set mu plus}
        \mu_+ = \lambda_k^{(j_k)}+\frac{\sqrt{M}}{3}\|\nabla f(y_k+d_k^{(j_k-1)})\|^{1/2}> \frac{2\sqrt{M}}{3}\|\nabla f(y_k+d_k^{(j_k-1)})\|^{1/2}\\
    \label{eq.newton equation}
        \left (\nabla ^2 f(y_k)+\mu_+ I\right) d_k^{(j_k)}=-\nabla f(y_k)\\
    \label{eq.step size}
        \|d_k^{(j_k)}\| = \frac{1}{3\sqrt{M}}\|\nabla f(y_k+d_k^{(j_k-1)})\|^{1/2}.
    \end{gather}
For simplicity, we omit the subscript $k$ and superscript $j_k$ in the following analysis, remember that 
\begin{equation*}
    y= y_k, \ x = y_k+d_k^{(j_k-1)}, \ \lambda = \lambda_k^{(j_k)}, \ d=d_k^{(j_k)}.
\end{equation*}

We first define the auxiliary function $g_y(\mu) = \frac{\mu}{\|\left (\nabla^2 f(y)+\mu I\right)^{-1}\nabla f(y)\|}$, given current point $y$, our goal is to find the regularize $\mu$ such that 
\begin{equation*}
    M \leq g_y(\mu) \leq 2M.
\end{equation*}
We first analyze some basic properties of $g_y(\mu)$,
\begin{lemma}
    \label{lem.property g}
    For any $y\in \mathbb{R}^n$, $g_y(\mu)$ is continuously differentiable and monotonically increasing for $\mu \in \left (0,+\infty\right ]$ and
    \begin{equation}
        \label{eq.derivative}
        g_y'(\mu) = \frac{1}{\|\left (\nabla ^2 f(y)+\mu I \right)^{-1}\nabla f(y)\|}+\frac{\mu}{\|\left (\nabla ^2 f(y)+\mu I \right)^{-1}\nabla f(y)\|^3}\sum_{i=1}^n\frac{\beta_i^2}{\left (\lambda_i+\mu\right )^3},
    \end{equation}
    and it is bounded both below and above
    \begin{equation}
        \label{eq.derivative bound}
        0< g_y'(\mu) \leq \frac{2}{\|\left (\nabla ^2 f(y)+\mu I \right)^{-1}\nabla f(y)\|}.
    \end{equation}
    Where $\lambda_i$ is the $i$-th eigenvalue of $\nabla ^2 f(y)$, $\beta_i = \nabla f(y)^T v_i$, $v_i$ denotes the eigenvector corresponding to $\lambda_i$.
\end{lemma}
\begin{proof}
    Since $\nabla^2 f(y)+\mu I \succ 0$ for $\mu>0$, we can apply eigen decomposition to it, then 
    \begin{equation*}
        \|\left (\nabla ^2 f(y)+\mu I \right)^{-1}\nabla f(y)\|=\sqrt{\sum_{i=1}^n \frac{\beta_i^2}{\left (\lambda_i+\mu\right )^2}}, \ g_y(\mu) = \frac{\mu}{\sqrt{\sum_{i=1}^n \frac{\beta_i^2}{\left (\lambda_i+\mu\right )^2}}}.
    \end{equation*}
    Obviously $g_y(\mu)$ is increasing in $\mu \in \left (0,+\infty\right ]$. Also, through some basic calculations, we have
    \begin{equation*}
        \begin{aligned}
            g_y'(\mu) &= \frac{1}{\sqrt{\sum_{i=1}^n \frac{\beta_i^2}{\left (\lambda_i+\mu\right )^2}}}+\mu \frac{1}{\left ( (\sum_{i=1}^n \frac{\beta_i^2}{\left (\lambda_i+\mu\right )^2}\right )^{3/2}}\sum_{i=1}^n\frac{\beta_i^2}{\left (\lambda_i+\mu\right )^3}\\
            &= \frac{1}{\|\left (\nabla ^2 f(y)+\mu I \right)^{-1}\nabla f(y)\|}+\frac{\mu}{\|\left (\nabla ^2 f(y)+\mu I \right)^{-1}\nabla f(y)\|^3}\sum_{i=1}^n\frac{\beta_i^2}{\left (\lambda_i+\mu\right )^3}\\
            &\leq \frac{1}{\|\left (\nabla ^2 f(y)+\mu I \right)^{-1}\nabla f(y)\|}+\frac{\mu}{\|\left (\nabla ^2 f(y)+\mu I \right)^{-1}\nabla f(y)\|^3}\sum_{i=1}^n\frac{\beta_i^2}{\mu\left (\lambda_i+\mu\right )^2}\\
            &=  \frac{2}{\|\left (\nabla ^2 f(y)+\mu I \right)^{-1}\nabla f(y)\|}.
        \end{aligned}
    \end{equation*}
    The inequality is because of $\lambda_i \geq 0$ for $i=1,\ldots,n$.
\end{proof}
In \autoref{alg.accelerated utr}, we already have the bracket points for the bisection, we now show that the choices are valid for the bisection procedure,
\begin{lemma}
    \label{eq.valid bracketing}
    As in the $k$-th iteration of \autoref{alg.accelerated utr}, we set
    \begin{equation*}
    \mu_- = \frac{\sqrt{M}}{3}\|\nabla f(x)\|^{1/2}, \mu_+ = \lambda+\frac{\sqrt{M}}{3}\|\nabla f(x)\|^{1/2},
    \end{equation*}
    then 
    \begin{equation*}
        g_y(\mu_-) < M, \ g_y(\mu_+) >2M.
    \end{equation*}
\end{lemma}
\begin{proof}
    From the mechanism of \autoref{alg.accelerated utr}, we have 
    \begin{equation*}
        \| \left (\nabla ^2 f(y)+\mu^+ I\right )^{-1} \nabla f(y)\| = \frac{1}{3\sqrt{M}}\|\nabla f(x)\|^{1/2}, \ \mu_+ >\frac{2\sqrt{M}}{3}\|\nabla f(x)\|^{1/2}=2M\|d\|,
    \end{equation*}
    thus $g_y(\mu_+) >2M$. Since $\mu_- <\mu^+$, then
    \begin{equation*}
        \| \left (\nabla ^2 f(y)+\mu^- I\right )^{-1} \nabla f(y)\|>\| \left (\nabla ^2 f(y)+\mu^+ I\right )^{-1} \nabla f(y)\| = \frac{1}{3\sqrt{M}}\|\nabla f(x)\|^{1/2},
    \end{equation*}
    therefore $g_y(\mu_-)< M$.
\end{proof}
From the continuity and monotonicity of $g_y(\mu)$, we know there exist $\mu_-<\mu_l<\mu_u<\mu_+$ such that
\begin{equation}\label{eq.target interval bracketing points}
    g_{y}(\mu_l)=M, \ g_{y}(\mu_u)=2M.
\end{equation}
To clarify, now we are given the initial interval $[\mu_-,\mu_+]$ and are using bisection to find some $\mu\in[\mu_l,\mu_u]$, once we establish a lower bound for the length of the target interval then we can bound the complexity for the bisection procedure.
\begin{lemma}
    \label{lem.bound target interval}
    For $\mu_l,\mu_u$ in \eqref{eq.target interval bracketing points}, we have
    \begin{equation}
        \label{eq.bound target interval}
        \mu_u-\mu_l \geq \frac{\sqrt{M}\|\nabla f(x)\|^{1/2}}{6}.
    \end{equation}
\end{lemma}
\begin{proof}
    By the mean value theorem, there exists $\xi\in \left [\mu_l,\mu_u\right]$, such that
    \begin{equation*}
        M=g_y(\mu_u)-g_y(\mu_l) = g_y'(\xi)(\mu_u-\mu_l),
    \end{equation*}
    combine the above with \eqref{eq.derivative bound} and \eqref{eq.target interval bracketing points} we have
    \begin{equation*}
        \begin{aligned}
            \mu_+-\mu_- &=  \frac{1}{g_y'(\xi)}\left(g_y(\mu_u)-g_y(\mu_l)\right)\\
            &\geq \frac{M\|\left (\nabla ^2 f(y)+\xi I \right )^{-1}\nabla f(y)\|}{2}\\
            &\geq \frac{M\|\left (\nabla ^2 f(y)+\mu_+ I \right )^{-1}\nabla f(y)\|}{2}\\
            &= \frac{\sqrt{M}\|\nabla f(x)\|^{1/2}}{6}.
        \end{aligned}
    \end{equation*}
    The third line is because of $\mu_+ >\xi$, the last line is from \eqref{eq.newton equation} and \eqref{eq.step size}.
\end{proof}
Now we established a lower bound for the target interval, it remains to bound the length of the initial interval $\mu_+ -\mu_- = \lambda$.
\begin{lemma}
    \label{lem.bound initial interval}
    At the $k$-th iteration of \autoref{alg.accelerated utr}, if the bisection begins, we have 
    \begin{equation*}
    \lambda\leq \frac{3\sqrt{M}\|\nabla f(y)\|}{\|\nabla f(x)\|^{1/2}}.
    \end{equation*}
    As a consequence, the complexity of the bisection search procedure is 
    \begin{equation}
        \label{eq.complexity bisection}
        O \left (\log \left( \frac{\mu_+-\mu_-}{\mu_u-\mu_l}\right)\right)= O\left (\log \left(\frac{\|\nabla f(y)\|}{\|\nabla f(x)\|}\right)\right)
    \end{equation}
\end{lemma}
\begin{proof}
    Note that by the mechanism of \autoref{alg.accelerated utr}, when bisection begins, we have
    \begin{equation*}
        \|\left (\nabla^2 f(y)+\mu_+ I \right)d\| =\|\nabla f(y)\|, \ \|d\| =\frac{1}{3\sqrt{M}}\|\nabla f(x)\|^{1/2}.
    \end{equation*}
    Note that $\nabla ^2 f(y)\succeq 0$, and $\mu_+ \geq \lambda$, we have
    \begin{equation*}
        \lambda\|d\| \leq \|\nabla f(y)\|,
    \end{equation*}
    thus 
    \begin{equation*}
        \lambda \leq \frac{\|\nabla f(y)\|}{\|d\|} = \frac{3\sqrt{M}\|\nabla f(y)\|}{\|\nabla f(x)\|^{1/2}}.
    \end{equation*}
    Combine the above with \eqref{eq.bound target interval} we have
    \begin{equation*}
        \frac{\mu_+-\mu_-}{\mu_u-\mu_l} \leq \frac{18\|\nabla f(y)\|}{\|\nabla f(x)\|},
    \end{equation*}
    hence we finished the proof.
\end{proof}
\subsection{The Total Complexity of \autoref{alg.accelerated utr}}
Note that as discussed in \autoref{coro.inner loop terminate}, when $k\in \mathcal{Z}$, there evolves a search procedure, whose complexity is $O(\log \frac{\|\nabla f(y_k)\|}{\epsilon})$, thus we need to establish an upper bound for $\| \nabla f(y_k)\|, k\in \mathcal{Z}$.
\begin{lemma}
    \label{lem.upper bound for zero lambda}
    Suppose \autoref{assm.bounded} and \autoref{assm.lipschitz} hold, for $k\in \mathcal{Z}$, we have 
    \begin{equation}
        \label{eq.upper bound for zero lambda}
        \|\nabla f(y_k)\|\leq \frac{9\kappa_H^2}{64M}.
    \end{equation}
\end{lemma}
\begin{proof}
    Note that when $k\in \mathcal{Z}$, we have $\lambda_k =0$, hence the following hold
    \begin{equation*}
        \left (\nabla^2 f(y_k)+\frac{\sqrt{M}}{3}\|\nabla f(y_k)\|^{1/2} I\right )d_k^{(0)} = -\nabla f(y_k), \ \|d_k^{(0)}\| \leq \frac{1}{3\sqrt{M}}\|\nabla f(y_k)\|^{1/2}.
    \end{equation*}
    Therefore, by triangle inequality, it holds that
    \begin{equation*}
        \|\nabla f(y_k)\| \leq \|\nabla^2 f(y_k) d_k^{(0)}\| +\frac{\sqrt{M}}{3}\|\nabla f(y_k)\|^{1/2}\|d_k^{(0)}\| \leq \|\nabla^2 f(y_k) d_k^{(0)}\|+\frac{1}{9}\|\nabla f(y_k)\|.
    \end{equation*}
    By \autoref{assm.bounded}, we have 
    \begin{equation*}
        \begin{aligned}
            \kappa_H \cdot \frac{1}{3\sqrt{M}}\|\nabla f(y_k)\|^{1/2} &\geq \kappa_H \|d_k^{(0)}\| \\
            & \geq \|\nabla^2 f(y_k) d_k^{(0)}\|\\
            & \geq \frac{8}{9}\|\nabla f(y_k)\|.
        \end{aligned}
    \end{equation*}
    Rearrange items we will have \eqref{eq.upper bound for zero lambda}.
\end{proof}
Now we can present the total complexity of \autoref{alg.accelerated utr}, 
\begin{theorem}
    \label{thm.total complexity newton}
    Suppose \autoref{assm.bounded} and \autoref{assm.lipschitz} hold, it takes at most 
    \begin{equation}\label{eq.total complexity newton}
        O\left (\epsilon^{-1/3}\log \left(\kappa_H/M\epsilon\right)\right)
    \end{equation}
    UTR or linear system oracles for \autoref{alg.accelerated utr} to find a solution $x$ satisfies \eqref{eq.convex approximate solution}. 
\end{theorem}
\begin{proof}
    Note that the UTR oracle is called only once in each outer loop and each inner loop, and the linear system oracle is called only in \autoref{alg.bisection}. The complexity of the inner loop is bounded in \autoref{coro.inner loop terminate} and \autoref{lem.upper bound for zero lambda}. Also, the complexity of the bisection procedure is bounded in \autoref{lem.bound initial interval}. Hence the total number of calls of such two oracles is bounded by \eqref{eq.total complexity newton}.
\end{proof}
\section{The Monteiro-Svaiter Acceleration}
Now we consider another variant of the accelerated Universal Trust-Region method, which utilized the framework proposed by \citet{monteiro2013accelerated}, here we make the additional assumption:
\begin{assumption}
    \label{assm.solvable}
    The problem \eqref{eq.main prob} is solvable, i.e., there exist $x^*$ such that 
    \begin{equation*}
        f(x^*) = f^* = \min_{x\in\mathbb{R}^n} f(x),
    \end{equation*}
    denote $D_0 = \|x_0-x^*\|.$
\end{assumption}

We first give the main \autoref{alg.ms accelerated utr} and its auxiliary search procedure \autoref{alg.bisection ms} as below:

\begin{algorithm}[ht]
    \caption{Accelerated Trust-Region Method via MS Framework(ATR\_MS)}\label{alg.ms accelerated utr}
    \begin{algorithmic}[1]
        \STATE{\textbf{input:} $x_0=v_0\in \mathbb{R}^n$, \ $A_0 =0$, \ $\theta>1$, $\gamma \leq \frac{1}{\theta}$, \ $0<\eta<1$}
        \FOR{$k=0, 1, \ldots, \infty$}
        \STATE{$y_k,d_k,a_k,\lambda_k,\sigma_k = \operatorname{BISECTION\_ATR\_MS}(x_k,v_k,A_k,\eta,\theta)$}
        \STATE{Update
        \begin{gather}
        \label{eq.update x}
            x_{k+1} = y_k+d_k\\
            \label{eq.update v}
            v_{k+1} = v_k-\gamma a_k \nabla f(x_{k+1})\\
            \label{eq.update A}
            A_{k+1} = A_k+a_k
        \end{gather}
        }
        \ENDFOR
    \end{algorithmic}
\end{algorithm}
\begin{algorithm}[!ht]
    \caption{Bisection for AT\_RMS}\label{alg.bisection ms}
    \begin{algorithmic}[1]
        \STATE{\textbf{input:} $x,v\in \mathbb{R}^n$, $A,\eta,\theta \in \mathbb{R}$;}
        \STATE{Define the curve:
        \begin{gather}
        \label{eq.update yk}
        y(\sigma) = \frac{A}{A+a(\sigma)}x+\frac{a(\sigma)}{A+a(\sigma)}v\\
        \label{eq.update ak}
        a(\sigma) =\frac{1+\sqrt{1+2A\sigma}}{\sigma}
        \end{gather}
        }
        \STATE{Set $\sigma_- = \sqrt{\frac{2M\epsilon}{1+2\theta}}, \ \sigma_+=\sqrt{\frac{MG_0}{\eta}}$, $(\lambda_-,d_-) = \operatorname{UTR}(y(\sigma_-),\sigma_-,\frac{1}{M}\sigma_-)$}
        \IF{$\lambda_- \leq (\theta-1)\sigma_-$}
        \STATE{\textbf{Output $y = y(\sigma_-)+d_-$ }as an approximate solution}\label{line.early termination ms}
        \ELSE
        \WHILE{True}
        \STATE{
        \begin{gather}
        \notag
        \sigma= \frac{\sigma_-+\sigma_+}{2}\\
        \label{eq.ms utr oracle}
        (\lambda,d) = \operatorname{UTR}(y(\sigma),\sigma,\frac{1}{M}\sigma)
        \end{gather}
        }
        \IF{$\lambda = 0$ and $\|d\| <\frac{\eta}{M}\sigma$}\label{line.target bisection 1}
        \STATE{$\sigma_+=\sigma$}
        \ELSIF{$\lambda > (\theta-1)\sigma$}\label{line.target bisection 2}
        \STATE{$\sigma_- = \sigma$}
        \ELSE
        \STATE{\textbf{output} $y(\sigma),d,a(\sigma),\lambda,\sigma$}
        \ENDIF
        \ENDWHILE
        \ENDIF
    \end{algorithmic}
\end{algorithm}
We remark here that from line \ref{line.target bisection 1} and line \ref{line.target bisection 2} of \autoref{alg.bisection ms}, the output of the bisection satisfies 
\begin{equation}
    \label{eq.relation lambda sigma}
        0\leq \lambda_k\leq (\theta-1)\sigma_k, \ \|d_k\| \geq \frac{\eta}{M}\sigma_k,
\end{equation}
\eqref{eq.relation lambda sigma} can be implied by 
\begin{equation}
    \label{eq.stonger relation lambda sigma}
    0< \lambda_k \leq (\theta-1)\sigma_k,
\end{equation}
since by \eqref{eq.optcond coml slack} we have $\|d_k\| =\frac{1}{M}\sigma_k$ and thus \eqref{eq.relation lambda sigma} holds. 

Also, we make the following clarification, note that there exists an early termination at line \ref{line.early termination ms} in \autoref{alg.bisection ms}, this is because of the following observation:
\begin{lemma}
    \label{lem.early termination ms}
    Suppose \autoref{alg.bisection ms} terminates at line \ref{line.early termination ms}, denote $y_- = y(\sigma_-)+d_-$, we have
    \begin{equation*}
        \|\nabla f(y_-) \|\leq \epsilon.
    \end{equation*}
\end{lemma}
\begin{proof}
    From the Lipschitz continuity of the Hessian, we have
    \begin{equation*}
        \|\nabla f(y_-)-\nabla f(y(\sigma))-\nabla^2 f(y(\sigma))d_-\|\leq \frac{M}{2}\|d_-\|^2.
    \end{equation*}
    By triangle inequality and \eqref{eq.optcond firstorder}, we have
    \begin{equation*}
        \begin{aligned}
            \|\nabla f(y_-)\| &\leq \frac{M}{2}\|d_-\|^2+\| \nabla f(y(\sigma))+\nabla^2 f(y(\sigma))d_-\|\\ 
            &= \frac{M}{2}\|d_-\|^2+(\sigma_-+\lambda_-) \|d_-\|\\
            &\leq \frac{1+2\theta}{2M}\sigma_-^2\\
            &\leq \epsilon.
        \end{aligned}
    \end{equation*}
\end{proof}
\subsection{Global Convergence of \autoref{alg.ms accelerated utr}}
We first assume the bisection procedure is effective and we have \eqref{eq.relation lambda sigma} at hand, we analyze the complexity of the framework.

\begin{lemma}\label{lem.ms estimate}
In each iteration $k$ of \autoref{alg.ms accelerated utr}, the following holds:
    \begin{equation}
        \label{eq.helper bounded x}
        \frac{1}{2}\|v_{k+1}-x^* \|^2 +\gamma A_{k+1}\left (f(x_{k+1}-f^*)\right)+ \frac{3\gamma A_{k+1}\sigma_k}{4}\|d_k\|^2 \leq \frac{1}{2}\|v_k-x^*\|^2 +\gamma A_k \left (f(x_k)-f^*\right).
    \end{equation}
    % \begin{equation}\label{eq.ms estimate}
    %     \gamma A_{k+1}\left (f(x_{k+1})-f^* \right) + \frac{1}{2}\|v_{k+1}-x^*\|^2 +\frac{3\gamma M A_{k+1}}{4}\|d_k\|^3 \leq \gamma A_k \left (f(x_k)-f^*\right ) +\frac{1}{2}\|v_k-x^*\|^2,
    % \end{equation}
    which further implies
    \begin{equation}
        \label{eq.ms estimate advanced}
        \gamma A_{k+1}\left (f(x_{k+1})-f^* \right) + \frac{1}{2}\|v_{k+1}-x^*\|^2 +B_{k+1} \leq \frac{1}{2}\|v_0-x^*\|^2,
    \end{equation}
    where $B_{k+1} = \frac{3\gamma M}{4}\sum_{i=0}^{k}A_{i+1}\|d_i\|^3.$
\end{lemma}
\begin{proof}
    \begin{equation*}
    \footnotesize
        \begin{aligned}
            &\|v_{k+1}-x^*\|^2\\
            &= \|v_k-x^* -\gamma a_k \nabla f(x_{k+1})\|^2 \\
            &= \|v_k-x^*\|^2 +\gamma^2 a_k^2 \|\nabla f(x_{k+1})\|^2 -2\gamma a_k \langle \nabla f(x_{k+1}),v_k-x^*\rangle\\
            &= \|v_k-x^*\|^2 +\gamma^2 a_k^2 \|\nabla f(x_{k+1})\|^2 -2\gamma \langle \nabla f(x_{k+1}), \left (A_k+a_k\right)y_k-A_k x_k-a_k x^* \rangle\\
            &=\|v_k-x^*\|^2 +\gamma^2 a_k^2 \|\nabla f(x_{k+1})\|^2-2\gamma \langle \nabla f(x_{k+1}), \left (A_k+a_k\right)y_k-\left(A_k+a_k\right)x_{k+1}+A_k\left(x_{k+1}- x_k\right)+a_k \left(x_{k+1}-x^*\right) \rangle\\
            &\leq\|v_k-x^*\|^2 +\gamma^2 a_k^2 \|\nabla f(x_{k+1})\|^2-2\gamma \left (A_k+a_k\right)\langle \nabla f(x_{k+1}), y_k-x_{k+1} \rangle-2\gamma A_k \left (f(x_{k+1})-f(x_k)\right)-2\gamma a_k\left (f(x_{k+1})-f^*\right).
        \end{aligned}
    \end{equation*}
    The third line comes from \eqref{eq.update yk}, therefore we have
    \begin{equation*}
    \begin{aligned}
        &\|v_{k+1}-x^*\|^2 +2\gamma A_{k+1} \left (f(x_{k+1})-f^*\right)\\
        &\leq \|v_k-x^*\|^2+2\gamma A_k \left (f(x_k)-f^*\right )+\gamma^2 a_k^2 \|\nabla f(x_{k+1})\|^2-2\gamma \left (A_k+a_k\right)\langle \nabla f(x_{k+1}), y_k-x_{k+1} \rangle\\
        &\leq\|v_k-x^*\|^2+2\gamma A_k \left (f(x_k)-f^*\right )+\gamma^2 a_k^2 \|\nabla f(x_{k+1})\|^2-2\gamma \left (A_k+a_k\right)\left (\frac{\|\nabla f(x_{k+1})\|^2}{\sigma_k+\lambda_k}+\frac{3(\sigma_k+\lambda_k)}{4}\|d_k\|^2\right)\\
        &\leq \|v_k-x^*\|^2+2\gamma A_k \left (f(x_k)-f^*\right )+\left(\gamma^2 a_k^2-\frac{2\gamma(A_k+a_k)}{\theta \sigma_k}\right)\|\nabla f(x_{k+1})\|^2 -\frac{3\gamma(A_k+a_k)(\sigma_k+\lambda_k)}{2}\|d_k\|^2.
    \end{aligned}
    \end{equation*}
    Note that $\gamma \leq \frac{1}{\theta}$ and the way we update $a_k$ in \eqref{eq.update ak}, rearranging items gives \eqref{eq.helper bounded x}.
    % \begin{equation}
    %     \label{eq.helper bounded x}
    %     \frac{1}{2}\|v_{k+1}-x^* \|^2 +\gamma A_{k+1}\left (f(x_{k+1}-f^*)\right)+ \frac{3\gamma A_{k+1}\sigma_k}{4}\|d_k\|^2 \leq \frac{1}{2}\|v_k-x^*\|^2 +\gamma A_k \left (f(x_k)-f^*\right).
    % \end{equation}
    Note that $\sigma_k \geq M\|d_k\|$,
    \begin{equation*}
        \frac{1}{2}\|v_{k+1}-x^* \|^2 +\gamma A_{k+1}\left (f(x_{k+1}-f^*)\right)+ \frac{3\gamma M A_{k+1}}{4}\|d_k\|^3 \leq \frac{1}{2}\|v_k-x^*\|^2 +\gamma A_k \left (f(x_k)-f^*\right).
    \end{equation*}
    Summing up the above for $i=0$ to $i=k$ gives \eqref{eq.ms estimate advanced}.
\end{proof}
\begin{lemma}
    \label{lem.bounded iterate}
    For every $k\geq 0$, 
    \begin{equation}
        \label{eq.bounded iterate}
        \|x_k- x^*\| \leq \left(\frac{2}{\sqrt{3\gamma}}+1\right)D_0,
    \end{equation}
    as a consequence,
    \begin{equation}
        \label{eq.bounded iterate y}
        \|y_k(\sigma)-x^*\| \leq \left(\frac{2}{\sqrt{3\gamma}}+1\right)D_0, \ \forall \sigma>0.
    \end{equation}
\end{lemma}
\begin{proof}
    First, note that
    \begin{equation*}
    \begin{aligned}
        \|x_{k+1}-x^*\| &\leq \|x_{k+1}-y_k\| + \|y_k-x^*\|\\
        & \leq \|d_k\|+ \frac{A_k}{A_{k+1}} \|x_k-x^*\| +\frac{a_k}{A_{k+1}}\|v_k-x^*\| \\
        &\leq \frac{1}{A_{k+1}}\left(A_{k+1}\|d_k\| +A_k\|x_k-x^*\| +  a_k D_0\right)\\
        &\leq \frac{1}{A_{k+1}}\left (\sum_{i=0}^k A_{i+1}\|d_i\|\right)+D_0.
        \end{aligned}
    \end{equation*}
    The second line is from \eqref{eq.update yk}, the third line is from \eqref{eq.ms estimate advanced}, and the fourth line is derived by iterating 
    \begin{equation*}
        A_{k+1} \|x_{k+1}-x^*\|\leq A_{k+1}\|d_k\|+A_k\|x_k-x^*\|+a_k D_0,
    \end{equation*}
    which is derived by multiplying both sides of the third line by $A_{k+1}$.

    Summing up \eqref{eq.helper bounded x}, we have
    \begin{equation*}
        \frac{3\gamma}{4}\sum_{i=0}^k A_{i+1}\sigma_i \|d_i\|^2 \leq \frac{1}{2}D_0^2,
    \end{equation*}
    to bound $\sum_{i=0}^k A_{i+1}\|d_i\|$, we come to the optimization problem:
    \begin{equation*}
        \max _{\zeta \in \mathbb{R}^{k+1}_+}\left\{\sum_{i=0}^{k} A_{i+1}\zeta_i: \quad \sum_{i=0}^{k} A_{i+1} \sigma_i\zeta_i^2 \leq \frac{2}{3\gamma}D_0^2\right\},
    \end{equation*}
    through similar analysis in \citet[Lemma A.2]{monteiro2013accelerated} we have
    \begin{equation*}
        \sum_{i=0}^k A_{i+1}\|d_i\| \leq \sqrt{\frac{2}{3\gamma}}\cdot \sqrt{\sum_{i=0}^k \frac{A_{i+1}}{\sigma_i}} D_0,
    \end{equation*}
    therefore,
    \begin{equation}\label{eq.bound initial x}
        \begin{aligned}
             \|x_{k+1}-x^*\| &\leq \frac{1}{A_{k+1}}\left (\sum_{i=0}^k A_{i+1}\|d_i\|\right)+D_0\\
             &\leq \left(\sqrt{\frac{2}{3\gamma}}\cdot \sum_{i=0}^k\sqrt{\frac{1}{\sigma_i}}\cdot \sqrt{\frac{1}{A_{k+1}}} +1\right)D_0.
        \end{aligned}
    \end{equation}
    The second line is from the fact that $A_k$ is monotone and $2$-norm is majorized by $1$-norm. From \eqref{eq.update ak} we have
    \begin{equation*}
        a_k \geq \frac{1}{\sigma_k}+\sqrt{\frac{2A_k}{\sigma_k}},
    \end{equation*}
    hence
    \begin{equation*}
        A_{k+1}\geq A_k+\frac{1}{2\sigma_k}+\sqrt{\frac{2A_k}{\sigma_k}},
    \end{equation*}
    taking the square root of both sides,
    \begin{equation*}
        \sqrt{A_{k+1}}\geq \sqrt{A_k}+\sqrt{\frac{1}{2\sigma_k}},
    \end{equation*}
    summing up from $i=0$ to $k$,
    \begin{equation*}
        \sqrt{A_{k+1}}\geq \sum_{i=0}^k \sqrt{\frac{1}{2\sigma_i}}.
    \end{equation*}
    Plugging the above into \eqref{eq.bound initial x}, we have
    \begin{equation*}
        \|x_{k+1}-x^*\| \leq \left( \frac{2}{\sqrt{3\gamma}}+1\right) D_0.
    \end{equation*}
    To prove the boundedness of $y_k(\sigma)$, just note that $y_k(\sigma)$ is a linear combination of $v_k$ and $x_k$, and $\|v_k-x^*\| \leq D_0$ as in \eqref{eq.ms estimate advanced}.
\end{proof}
Now with the help of \autoref{lem.ms estimate}, it only requires analyzing the growth rate of $A_k$ to derive the convergence result of \autoref{alg.ms accelerated utr}, the growth rate has already been studied in \citet[Lemma 4.2]{monteiro2013accelerated}, \citet[Lemma 4.3.5]{nesterov_lectures_2018}, we defer the proof to the appendix.
\begin{lemma}
    \label{lem.growth of Ak}
    For and $k\geq 1$, we have
    \begin{equation}
        \label{eq.growth of Ak}
        A_k \geq \left(\frac{\eta}{2}\left (\frac{3\gamma}{2M^2D_0^2}\right)^{1/3}\right)^{3 / 2}\left(\frac{2 k+1}{3}\right)^{7/2} = \Omega(k^{7/2}).
    \end{equation}
\end{lemma}
Hence the convergence of \autoref{alg.ms accelerated utr} comes as follows:
\begin{theorem}\label{thm.outer loop complexity}
    Suppose in each iteration $k$ the bisection is effective, then for any $0<\epsilon<1$, \autoref{alg.ms accelerated utr} find an $\epsilon$-approximate solution as in \eqref{eq.convex approximate solution} in $O(\epsilon^{-2/7})$ iterations.
\end{theorem}
\begin{proof}
    Note that suppose the bisection procedure is effective, then it either terminates at line \ref{line.early termination ms} or return some $\sigma_k$ satisfies \eqref{eq.relation lambda sigma}. Then from \eqref{eq.ms estimate advanced} and \eqref{eq.growth of Ak} we finished the proof.
\end{proof}
\subsection{Bisection Procedure}
Next, we elaborate on how \autoref{alg.bisection ms} is guaranteed to find $\sigma_k$ with \eqref{eq.relation lambda sigma} holds. First, we introduce several auxiliary results to help the analysis.
\begin{corollary}
    \label{coro.boundedness ysigma}
    For any $k$ and any $\sigma$, 
    \begin{equation}
        \label{eq.ub yk sigma}
        \|\nabla f(y_k(\sigma))\| \leq G_0:= \left(\frac{2}{\sqrt{3\gamma}}+1\right)\|\nabla ^2 f(x^*)\|D_0+\frac{M}{2}\left(\frac{2}{\sqrt{3\gamma}}+1\right)^2 D_0^2.
    \end{equation}
\end{corollary}
\begin{proof}
    We have
    \begin{equation*}
            \|\nabla f(y_k(\sigma)) - \nabla f(x^*) -\nabla^2 f(x^*)(y_k(\sigma)-x^*)\| \leq \frac{M}{2}\|y_k(\sigma)-x^*\|^2.
    \end{equation*}
    By triangle inequality, we have
    \begin{equation*}
        \begin{aligned}
            \|\nabla f(y_k(\sigma)) \| &\leq \|\nabla f(x^*)+\nabla^2 f(x^*)(y_k(\sigma)-x^*)\| +\frac{M}{2}\|y_k(\sigma)-x^*\|^2\\
            &\leq \|\nabla^2 f(x^*)(y_k(\sigma)-x^*)\|+\frac{M}{2}\|y_k(\sigma)-x^*\|^2\\
            &\leq \left(\frac{2}{\sqrt{3\gamma}}+1\right)\|\nabla ^2 f(x^*)\|D_0+\frac{M}{2}\left(\frac{2}{\sqrt{3\gamma}}+1\right)^2 D_0^2.
        \end{aligned}
    \end{equation*}
\end{proof}


We can define the quantity that is crucial in the analysis, here $y$ can be arbitrary and independent of $\sigma$:
\begin{equation}
    \label{eq.linesearch quantity}
    \psi_(\sigma,y) := \frac{1}{\sigma}\left \| \left(\nabla^2 f(y)+\sigma I \right)^{-1}\nabla f(y)\right \|, \ \sigma>0, y\in\mathbb{R}^n.
\end{equation}
We first analyze its value with respect to $\sigma$,
\begin{lemma}
    \label{lem.psi property wrt sigma}
    For any $y\in\mathbb{R}^n$, suppose $0<\sigma\leq \Bar{\sigma}$, then we have
    \begin{equation}
        \label{eq.psi property wrt sigma}
        \left(\frac{\sigma}{\Bar{\sigma}}\right)^2 \psi(\sigma,y)\leq \psi(\Bar{\sigma},y) \leq \frac{\sigma}{\Bar{\sigma}}\psi(\sigma,y).
    \end{equation}
\end{lemma}
\begin{proof}
    First note that 
    \begin{equation*}
        \sigma \psi(\sigma,y) = \left \| \left(\nabla^2 f(y)+\sigma I \right)^{-1}\nabla f(y)\right \|,
    \end{equation*}
    since $\nabla^2 f(y)\succeq 0$, we have
    \begin{equation*}
        \sigma \psi(\sigma,y) \geq \Bar{\sigma} \psi(\Bar{\sigma},y),
    \end{equation*}
    which is the second argument. Similarly,
    \begin{equation*}
        \sigma^2 \psi(\sigma,y) = \left \| \left(\frac{\nabla^2 f(y)}{\sigma}+ I \right)^{-1}\nabla f(y)\right \|,
    \end{equation*}
    then $\sigma^2 \psi(\sigma,y)\leq \Bar{\sigma}^2 \psi(\Bar{\sigma},y)$, which finished the proof.
\end{proof}
Next, we analyze the value of $\psi(\sigma,y)$ with respect to $y$, since its proof is quite lengthy and resembles the proof of \citet[Lemma 7.5]{monteiro2013accelerated}, we defer it to the appendix.
\begin{lemma}
    \label{lem.psi property wrt y}
    For any $y,\Bar{y} \in \mathbb{R}^n$ and $\sigma>0$, then 
    \begin{equation}
        \label{eq.psi property wrt y}
       \left \vert \psi(\sigma,y)-\psi(\sigma,\Bar{y}) \right\vert \leq \frac{1}{\sigma}\|y-\Bar{y}\| +\frac{M}{\sigma^2}\|y-\Bar{y}\|^2 +\frac{2M}{\sigma}\|y-\Bar{y}\|\delta,
    \end{equation}
    where $\delta:=\min\{\psi(\sigma,\Bar{y}),\psi(\sigma,y)\}$. Further, we have
    \begin{equation}
        \label{eq.psi property wrt y advanced}
        \psi(\sigma,y) \leq \frac{1}{\sigma}\|y-\Bar{y}\| +\frac{M}{\sigma^2}\|y-\Bar{y}\|^2 +\left(\frac{2M}{\sigma}\|y-\Bar{y}\|+1\right)\psi(\sigma,\Bar{y}).
    \end{equation}
\end{lemma}
\begin{lemma}
    \label{lem.reduce line search}
    In the bisection procedure, at line \ref{line.target bisection 1}, if  
    \begin{equation}
    \label{eq.bracket small lambda}
    \lambda = 0, \ \|d\|<\frac{\eta}{M}\sigma,
    \end{equation}
then $\psi(\sigma,y(\sigma))<\frac{\eta}{M}$.  Otherwise, if at line \ref{line.target bisection 2},
\begin{equation}
    \label{eq.bracket big lambda}
    \lambda > (\theta-1)\sigma,
\end{equation}
then $\psi(\sigma,y(\sigma))<\frac{1}{M}.$
\end{lemma}
\begin{proof}
If $\sigma$ belongs to case \eqref{eq.bracket small lambda}, we have 
\begin{equation*}
    \|d\| = \left \|\left(\nabla^2 f(y_k(\sigma))+\sigma I\right)^{-1}\nabla f(y(\sigma))\right\|<\frac{\eta}{M}\sigma,
\end{equation*}
which is $\psi(\sigma_k,y_k(\sigma_k))<\frac{\eta}{M}$.

Then consider the other case \eqref{eq.bracket big lambda}, we have
\begin{equation*}
    \|d\| = \left \|\left(\nabla^2 f(y(\sigma))+\sigma I+\lambda I\right)^{-1}\nabla f(y(\sigma))\right\|=\frac{1}{M}\sigma,
\end{equation*}
from \eqref{eq.psi property wrt sigma}, we have 
\begin{equation*}
     \left \|\left(\nabla^2 f(y(\sigma))+\sigma I\right)^{-1}\nabla f(y(\sigma))\right\|>\left \|\left(\nabla^2 f(y(\sigma))+\sigma I+\lambda I\right)^{-1}\nabla f(y(\sigma))\right\|=\frac{1}{M}\sigma,
\end{equation*}
which means $\psi \left(\sigma,y(\sigma)\right)>\frac{1}{M}$.
\end{proof}
\begin{corollary}
    \label{coro.reduce line search}
    In the bisection procedure, if $\sigma$ satisfies
    \begin{equation}
        \label{eq.reduce line search interval}
        \frac{\eta}{M}\leq \psi_(\sigma,y(\sigma)) \leq \frac{1}{M},
    \end{equation}
    then \eqref{eq.ms utr oracle} and \eqref{eq.relation lambda sigma} hold.
\end{corollary}
The above \autoref{coro.reduce line search} means we can reduce the line search problem to the one in \citet{monteiro2013accelerated}. We will provide a simplified analysis of our setting here. We still use the bisection method in our theoretical analysis and start with finding the bracketing point,
\begin{lemma}
    \label{lem.bracketing points}
    In \autoref{alg.bisection ms}, we let 
    \begin{equation}
        \label{eq.bracket points}
        \sigma_- = \sqrt{\frac{2M\epsilon}{1+2\theta}}, \ \sigma_+ = \sqrt{\frac{MG_0}{\eta}},
    \end{equation}
    then we either have
    \begin{equation*}
        \psi\left (\sigma_-,y(\sigma_-)\right) >\frac{1}{M}, \ \psi\left (\sigma_+,y(\sigma_+)\right) <\frac{\eta}{M}
    \end{equation*}
    or the algorithm early terminated at line \ref{line.early termination ms}.
\end{lemma}
\begin{proof}
    First, we show that $\psi(\sigma_+,y(\sigma_+))<\frac{\eta}{M}$,
    \begin{equation*}
        \begin{aligned}
        \psi(\sigma_+,y(\sigma_+)) &= \frac{1}{\sigma_+}\left \|\left(\nabla^2 f(y(\sigma_+))+\sigma_+ I\right)^{-1}\nabla f(y(\sigma_+))\right\| \\
        &\leq \frac{1}{\sigma_+^2}\|\nabla f(y(\sigma_+))\| \\
        &\leq \frac{1}{\sigma_+^2} G_0\\
        &= \frac{\eta}{M}.
        \end{aligned}
    \end{equation*}
    The second line is from $\nabla^2 f(y(\sigma_+))\succeq 0$, and the third line is from \eqref{eq.ub yk sigma}. 
    
    From \autoref{lem.early termination ms}, we can conclude that if the algorithm does not early terminate, we have $\lambda_- > (\theta-1)\sigma_-$, from \autoref{lem.reduce line search} we can conclude $\psi(\sigma_-,y(\sigma_-))>\frac{1}{M}.$
\end{proof}
After finding the bracket points, we are now ready to analyze the complexity of the bisection procedure, first, we need to introduce the following property of $y_k(\cdot)$, since it is an auxiliary lemma and the proof is almost the same as in \citet{monteiro2013accelerated}, we delay its proof to the appendix.
\begin{lemma}[Lemma 7.13, \citet{monteiro2013accelerated}]
\label{lem.curve condition}
    For each $k$, the curve $y_k(\cdot)$ satisfies
    \begin{equation}
        \label{eq.curve condition}
        \|y_k(s)-y_k(t)\| \leq \frac{M_0}{t}(s-t), \ \forall s\geq t >0,
    \end{equation}
    where $M_0 = 2\left (\frac{2}{\sqrt{3\gamma}}+1\right) D_0$.
\end{lemma}
Note that if the bisection procedure does not stop, we have
\begin{gather}
    \label{eq.bisection condition}
    \psi_+:=\psi(\sigma_+,y(\sigma_+))<\frac{\eta}{M}, \ \psi_-:=\psi(\sigma_-,y(\sigma_-))>\frac{\theta}{M},\\
    \label{eq.bracket points condition}
    \sigma_- \geq \sqrt{\frac{2M\epsilon}{1+2\theta}}, \ \sigma_+ \leq \sqrt{\frac{MG_0}{\eta}}.
\end{gather}
Besides, the gap between them has the following bound,
\begin{lemma}
    \label{lem.bound gap bisection}
    During the bisection procedure, we have
    \begin{equation}\label{eq.bound gap bisection}
    \begin{aligned}
    &\psi_- -\psi_+ \\
    &\leq \left(\frac{\sigma_+}{\sigma_-} \right)^2 \left[\frac{1}{\sigma_+}\|y(\sigma_-)-y(\sigma_+)\|+\frac{M}{\sigma_+^2}\|y(\sigma_-)-y(\sigma_+)\|^2+\frac{2M}{\sigma_+}\|y(\sigma_-)-y(\sigma_+)\|\psi_+\right]+\left[\left(\frac{\sigma_+}{\sigma_-}\right)^2-1\right]\psi_+.
    \end{aligned}
    \end{equation}
    Further, it gives
    \begin{equation}
        \label{eq.bounded gap bisection sigma}
        \begin{aligned}
             &\psi_- -\psi_+ \\
             &\leq \frac{\sigma_+}{\sigma_-^2}\left[\frac{M_0}{\sigma_-}+\frac{MM_0^2}{\sigma_-^2}+2\left(\frac{MM_0}{\sigma_-}+1\right)\psi_+
             \right](\sigma_+-\sigma_-).
        \end{aligned}
    \end{equation}
\end{lemma}
\begin{proof}
    \begin{equation*}
        \begin{aligned}
            \psi_- -\psi_+&= \psi(\sigma_-,y(\sigma_-))-\psi(\sigma_+,y(\sigma_+))\\
            &\leq \left(\frac{\sigma_+}{\sigma_-}\right)^2\psi(\sigma_+,y(\sigma_-))-\psi(\sigma_+,y(\sigma_+))\\
            &\leq \left(\frac{\sigma_+}{\sigma_-}\right)^2 \left[\frac{1}{\sigma_+}\|y(\sigma_-)-y(\sigma_+)\|+\frac{M}{\sigma_+^2}\|y(\sigma_-)-y(\sigma_+)\|^2+\left(\frac{2M}{\sigma_+}\|y(\sigma_-)-y(\sigma_+)\|+1\right)\psi_+\right]-\psi_+\\
            &=\left(\frac{\sigma_+}{\sigma_-} \right)^2 \left[\frac{1}{\sigma_+}\|y(\sigma_-)-y(\sigma_+)\|+\frac{M}{\sigma_+^2}\|y(\sigma_-)-y(\sigma_+)\|^2+\frac{2M}{\sigma_+}\|y(\sigma_-)-y(\sigma_+)\|\psi_+\right]+\left[\left(\frac{\sigma_+}{\sigma_-}\right)^2-1\right]\psi_+.
        \end{aligned}
    \end{equation*}
    The second line is because of \autoref{lem.psi property wrt sigma}, the third line is because of \autoref{lem.psi property wrt y}. To derive \eqref{eq.bounded gap bisection sigma}, just plug \eqref{eq.curve condition} into \eqref{eq.bound gap bisection}.
\end{proof}
The above lemma directly leads to the following:
\begin{corollary}
    \label{coro.bisection bound}
    During the bisection procedure, we have
    \begin{equation*}
    \begin{aligned}
        &\frac{\sigma_+}{\sigma_-^2}\left[\frac{M_0}{\sigma_-}+\frac{MM_0^2}{\sigma_-^2}+2\left(\frac{MM_0}{\sigma_-}+1\right)\psi_+
             \right]\\
             &\leq \sqrt{G_0}\left(\frac{(1+2\theta)^{3/2}M_0}{2\sqrt{2}M\eta \epsilon^{3/2}}+\frac{(1+2\theta)^2M_0^2}{4\sqrt{M}\eta\epsilon^2}+\frac{(1+2\theta)^{3/2}M_0}{\sqrt{2}M\epsilon^{3/2}}+\frac{(1+2\theta)\sqrt{G_0}}{M^{3/2}\epsilon}\right).
    \end{aligned}
    \end{equation*}
    Denote $C(G_0,M_0,M,\epsilon):= M\sqrt{G_0}\left(\frac{(1+2\theta)^{3/2}M_0}{2\sqrt{2}M\eta \epsilon^{3/2}}+\frac{(1+2\theta)^2M_0^2}{4\sqrt{M}\eta\epsilon^2}+\frac{(1+2\theta)^{3/2}M_0}{\sqrt{2}M\epsilon^{3/2}}+\frac{(1+2\theta)\sqrt{G_0}}{M^{3/2}\epsilon}\right)$                                     
                                , we have
    \begin{equation}
        \label{eq.bound value by brackets}
        \sigma_+-\sigma_->\frac{1-\eta}{C(G_0,M_0,M,\epsilon)}.
    \end{equation}
\end{corollary}
Now we are ready to analyze the total complexity of \autoref{alg.bisection ms}.
\begin{theorem}
    \label{thm.final complexity of msatr}
    Suppose \autoref{assm.lipschitz} holds, it takes at most 
    \begin{equation}
        \label{eq.final complexity of msatr}
        O\left(\epsilon^{-2/7}\log\left(MD_0/\epsilon\right)\right)
    \end{equation}
    UTR oracles for \autoref{alg.ms accelerated utr} to find a solution $x$ satisfies \eqref{eq.convex approximate solution}.
\end{theorem}
\begin{proof}
    We begin with analyzing the complexity of the bisection procedure, note that if the bisection procedure does not terminate, from \eqref{eq.bound value by brackets} we have
    \begin{equation*}
        \sigma_+-\sigma- > \frac{1-\eta}{ C(G_0,M_0,M,\epsilon)}.
    \end{equation*}
    Note that the bracket points are defined as in \eqref{eq.bracket points}, therefore by the mechanism of the bisection method, the total number of bisection in the $k$-th iteration $N_k$ is bounded by 
    \begin{equation*}
    \begin{aligned}
        N_k &\leq \log \left(\frac{\left(\sqrt{\frac{MG_0}{\eta}}-\sqrt{\frac{2M\epsilon}{1+2\theta}}\right)C(G_0,M_0,M,\epsilon)}{1-\eta}\right)\\
        &\leq \log \left(\frac{\sqrt{\frac{MG_0}{\eta}}C(G_0,M_0,M,\epsilon)}{1-\eta}\right).
    \end{aligned}
    \end{equation*}
    Since $G_0,M_0$ is polynomial in $D_0$, by the definition of $C(G_0,M_0,M,\epsilon)$ and omitting the algorithm parameters, we conclude $N_k \leq O\left(\log MD_0/\epsilon\right)$. From \autoref{thm.outer loop complexity} and the complexity of each iteration, we can conclude that \autoref{alg.ms accelerated utr} takes at most $O\left(\epsilon^{-2/7}\log \left(MD_0/\epsilon\right)\right)$ UTR oracles to find a point satisfying \eqref{eq.convex approximate solution}.
\end{proof}

\section{Numerical Experiments}
\clearpage
\addcontentsline{toc}{section}{References}

\bibliographystyle{unsrtnat}
\bibliography{ref}

\clearpage
\section*{Appendix}
\appendix
\subsection*{Proof to \autoref{lem.growth of Ak}}
\begin{proof}
    First note that
    \begin{equation*}
        A_{k+1}^{1/2}-A_k^{1/2} = \frac{a_k}{A_{k+1}^{1/2}+A_k^{1/2}}=\frac{1}{A_{k+1}^{1/2}+A_k^{1/2}}\sqrt{\frac{2A_{k+1}}{\sigma_k}}\geq \sqrt{\frac{1}{2\sigma_k}}.
    \end{equation*}
    Summing up the above from $i=0$ to $k-1$ gives 
    \begin{equation*}
        A_k \geq \frac{1}{2} \left (\sum_{i=0}^{k-1}\frac{1}{\sigma_i^{1/2}}\right)^2,
    \end{equation*}
    from \eqref{eq.relation lambda sigma}, the above gives
    \begin{equation}\label{eq.Ak helper 1}
        A_k \geq \frac{\eta}{2M} \left (\sum_{i=0}^{k-1} \frac{1}{\|d_i\|^{1/2}}\right)^2,
    \end{equation}
    on the other hand, from \eqref{eq.ms estimate advanced} we have
    \begin{equation*}
        B_k = \frac{3\gamma M}{4}\sum_{i=0}^{k-1} A_{i+1}\|d_i\|^3 \leq \frac{1}{2}\|v_0-x^*\|^2.
    \end{equation*}
    To estimate $A_k$ from below, define $\zeta_i = \|d_i\|^{1/2}$, $D = \frac{2}{3\gamma M}\|v_0-x^*\|^2$, we come to the following optimization problem
    \begin{equation*}
\zeta^*=\min _{\zeta \in \mathbb{R}^k}\left\{\sum_{i=0}^{k-1} \frac{1}{\zeta_i}: \quad \sum_{i=0}^{k-1} A_{i+1} \zeta_i^6 \leq D\right\}.
\end{equation*}
Introducing the Lagrangian multiplier $w$, the optimality condition gives 
\begin{equation*}
\frac{1}{\zeta_i^2}=w A_{i+1} \zeta_i^5, \quad i=0, \ldots, k-1
\end{equation*}
thus $\zeta_i = \left( \frac{1}{w A_{i+1}}\right)^{1/7}$, we have $w>0$ and the constraint is active, 
\begin{equation*}
D=\sum_{i=0}^{k-1} A_{i+1}\left(\frac{1}{w A_{i+1}}\right)^{6 / 7}=\frac{1}{w^{6 / 7}} \sum_{i=0}^{k-1} A_{i+1}^{1 / 7},
\end{equation*}
therefore $\zeta^*=\sum_{i=0}^{k-1}\left(w A_{i+1}\right)^{1 / 7}=\frac{1}{D^{1 / 6}}\left(\sum_{i=0}^{k-1} A_{i+1}^{1 / 7}\right)^{7 / 6}$, plugging back $\|d_i\| $ and $ \|v_0-x^*\|^2$, we have
\begin{equation*}
    \sum_{i=0}^{k-1}\frac{1}{\|d_i\|^{1/2}} \geq \left (\frac{3\gamma M}{2\|v_0-x^*\|^2}\right)^{1/6} \left(\sum_{i=0}^{k-1} A_{i+1}^{1 / 7}\right)^{7 / 6},
\end{equation*}
from \eqref{eq.Ak helper 1} we have
\begin{equation}
    \label{eq.Ak helper 2}
    A_k \geq \frac{\eta}{2M}\left (\frac{3\gamma M}{2\|v_0-x^*\|^2}\right)^{1/3} \left(\sum_{i=1}^{k} A_{i}^{1 / 7}\right)^{7 / 3}, \ k\geq 1.
\end{equation}
Denote $\omega = \frac{\eta}{2M}\left (\frac{3\gamma M}{2\|v_0-x^*\|^2}\right)^{1/3}$, $C_k =\left(\sum_{i=1}^k A_i^{1 / 7}\right)^{2 / 3}$, plugging them into \eqref{eq.Ak helper 2} we have
\begin{equation*}
    C_{k+1}^{3 / 2}-C_k^{3 / 2} \geq \omega^{1 / 7} C_{k+1}^{1 / 2},
\end{equation*}
which gives $C_1\geq \omega^{1/7}$ and 
\begin{equation*}
\begin{aligned}
\omega^{1 / 7} C_{k+1}^{1 / 2} & \leq\left(C_{k+1}^{1 / 2}-C_k^{1 / 2}\right)\left(C_{k+1}^{1 / 2}\left(C_{k+1}^{1 / 2}+C_k^{1 / 2}\right)+C_k\right) \\
& \leq\left(C_{k+1}^{1 / 2}-C_k^{1 / 2}\right)\left(C_{k+1}^{1 / 2}\left(C_{k+1}^{1 / 2}+C_k^{1 / 2}\right)+\frac{1}{2} C_{k+1}^{1 / 2}\left(C_{k+1}^{1 / 2}+C_k^{1 / 2}\right)\right) \\
& =\frac{3}{2} C_{k+1}^{1 / 2}\left(C_{k+1}-C_k\right).
\end{aligned}
\end{equation*}
Thus $C_k \geq \omega^{1 / 7}\left(1+\frac{2}{3}(k-1)\right), k \geq 1$. For $A_k$, by \eqref{eq.Ak helper 2} we have
\begin{equation*}
\begin{aligned}
A_k & \geq \omega\left(C_k^{3 / 2}\right)^{7 / 3} \geq \omega\left(\omega^{1 / 7} \cdot \frac{2 k+1}{3}\right)^{7 / 2}=\omega^{3 / 2}\left(\frac{2 k+1}{3}\right)^{7 / 2} \\
& =\left(\frac{\eta}{2}\left (\frac{3\gamma}{2M^2\|v_0-x^*\|^2}\right)^{1/3}\right)^{3 / 2}\left(\frac{2 k+1}{3}\right)^{3.5}
\end{aligned}
\end{equation*}
\end{proof}

\subsection*{Proof to \autoref{lem.psi property wrt y}}
\begin{proof}
    We denote 
    \begin{gather*}
        x = \arg\min_{x\in \mathbb{R}^n} \ \nabla f(y)^T(x-y) +\frac{1}{2}(x-y)\nabla ^2 f(y)(x-y) +\frac{\sigma}{2}\|x-y\|^2,\\
        \Bar{x} = \arg\min_{x\in \mathbb{R}^n} \ \nabla f(\Bar{y})^T(x-\Bar{y}) +\frac{1}{2}(x-\Bar{y})\nabla ^2 f(\Bar{y})(x-\Bar{y}) +\frac{\sigma}{2}\|x-\Bar{y}\|^2,
    \end{gather*}
    the optimality conditions of the above are 
    \begin{gather*}
        \nabla f(y) + \nabla ^2 f(y)(x-y) +\sigma(x-y)=0,\\
        \nabla f(\Bar{y})+\nabla ^2 f(\Bar{y})(\Bar{x}-\Bar{y})+\sigma(\Bar{x}-\Bar{y})=0.
    \end{gather*}
    Denote 
    \begin{gather*}
        v = \nabla f(y) + \nabla ^2 f(y)(x-y),\\
        \Bar{v} =  \nabla f(\Bar{y})+\nabla ^2 f(\Bar{y})(\Bar{x}-\Bar{y}),\\
        u = \nabla f(\Bar{y})+\nabla ^2 f(\Bar{y})(x-\Bar{y}).
    \end{gather*}
    Plugging them into the optimality conditions gives
    \begin{gather}
    \label{eq.def v}
        v+\sigma(x-y)=0,\\
    \label{eq.def bar v}
        \Bar{v}+\sigma(\Bar{x}-\Bar{y})=0,\\
    \label{eq.target distance initial}
        \left \vert \psi(\sigma,y)-\psi(\sigma,\Bar{y}) \right\vert= \left \vert\frac{1}{\sigma^2}\|v\| -\frac{1}{\sigma^2}\|\Bar{v}\|\right\vert\leq \frac{1}{\sigma^2}\|v-\Bar{v}\|.
    \end{gather}
    Subtracting \eqref{eq.def bar v} from \eqref{eq.def v}, we have
    \begin{equation*}
        \sigma(x-\Bar{x})+v -\Bar{v} = \sigma(y-\Bar{y}).
    \end{equation*}
    Plus both sides by $u$, we have
    \begin{equation*}
        \sigma(x-\Bar{x})+u-\Bar{v} = \sigma(y-\Bar{y})+u-v.
    \end{equation*}
    It is easy to show that
    \begin{equation*}
        \langle x-\Bar{x},u-\Bar{v}\rangle = \langle x-\Bar{x},\nabla ^2 f(\Bar{y})(x-\Bar{x})\rangle\geq 0,
    \end{equation*}
    thus from triangle inequality, we have
    \begin{equation*}
        \|u-\Bar{v}\| \leq \|\sigma(x-\Bar{x})+u-\Bar{v}\| \leq \sigma \|y-\Bar{y}\|+\|u-v\|,
    \end{equation*}
    hence 
    \begin{equation}\label{eq.bound psi 1}
        \|v-\Bar{v}\| \leq \sigma\|y-\Bar{y}\| +2\|u-v\|.
    \end{equation}
    Now we bound $\|u-v\|$,
    \begin{equation}\label{eq.bound psi 2}
    \begin{aligned}
        \|u-v\| &= \left\|\nabla f(\Bar{y})+\nabla ^2 f(\Bar{y})(x-\Bar{y})-\left(\nabla f(y) + \nabla ^2 f(y)(x-y)\right)\right\|\\
        &= \left\|\left(\nabla f(\Bar{y})+\nabla^2 f(\Bar{y})(y-\Bar{y})-\nabla f(y)\right)+\left(\nabla^2 f(\Bar{y})(x-y)-\nabla^2 f(y)(x-y)\right)\right\|\\
        &\leq \left\|\left(\nabla f(\Bar{y})+\nabla^2 f(\Bar{y})(y-\Bar{y})-\nabla f(y)\right)\right \| +\left \|\left(\nabla^2 f(\Bar{y})-\nabla^2 f(y)\right)\left(x-y\right)\right\|\\
        &\leq \frac{M}{2}\|y-\Bar{y}\|^2+M\|y-\Bar{y}\|\|x-y\|\\
        &=\frac{M}{2}\|y-\Bar{y}\|^2+\sigma M \|y-\Bar{y}\|\psi(\sigma,y).
    \end{aligned}
    \end{equation}
    Plugging \eqref{eq.bound psi 1} and \eqref{eq.bound psi 2} into \eqref{eq.target distance initial}, we have
    \begin{equation*}
        \left \vert \psi(\sigma,y)-\psi(\sigma,\Bar{y}) \right\vert \leq \frac{1}{\sigma}\|y-\Bar{y}\| +\frac{M}{\sigma^2}\|y-\Bar{y}\|^2 +\frac{2M}{\sigma}\|y-\Bar{y}\|\psi(\sigma,y).
    \end{equation*}
    Similarly, we can prove
    \begin{equation*}
        \left \vert \psi(\sigma,y)-\psi(\sigma,\Bar{y}) \right\vert \leq \frac{1}{\sigma}\|y-\Bar{y}\| +\frac{M}{\sigma^2}\|y-\Bar{y}\|^2 +\frac{2M}{\sigma}\|y-\Bar{y}\|\psi(\sigma,\Bar{y}).
    \end{equation*}
    Hence we have proved \eqref{eq.psi property wrt y}, by applying triangle inequality we have \eqref{eq.psi property wrt y advanced}.
\end{proof}
\subsection*{Proof to \autoref{lem.curve condition}}
\begin{proof}
    For any $k$, we have
    \begin{equation*}
    \begin{aligned}
        y_k(\sigma) &= \frac{A_k}{A_k+a_k(\sigma)}x_k+\frac{a_k(\sigma)}{A_k+a_k(\sigma)}v_k\\
        &= x_k + \tau (\sigma) (v_k-x_k),
    \end{aligned}
    \end{equation*}
    where $\tau(\sigma) = \frac{a_k(\sigma)}{A_k+a_k(\sigma)}.$ For any $s\geq t>0$,
    \begin{equation*}
        \|y_k(s)-y_k(t)\| = \left|\tau(s)-\tau(t)\right| \|v_k-x_k\|,
    \end{equation*}
    by the mean value theorem, we have
    \begin{equation*}
        \|y_k(s)-y_k(t)\| = \left|\tau'(\xi)\right|(s-t) \|v_k-x_k\|\leq M_0\left|\tau'(\xi)\right|(s-t),
    \end{equation*}
    where $\xi \in \left[t,s\right]$. Note that from \eqref{eq.update ak} we have
    \begin{equation*}
        \sigma a_k^2=2a_k+2A_k,
    \end{equation*}
    it leads to 
    \begin{equation*}
        \tau(\sigma) = \frac{a_k(\sigma)}{A_k+a_k(\sigma)} = \frac{2}{\sigma a_k(\sigma)} = \frac{2}{1+\sqrt{1+2A_k\sigma}}.
    \end{equation*}
    Its derivative is
    \begin{equation*}
        \tau'(\sigma) = -\frac{2A_k}{\sqrt{1+2A_k\sigma}\left(1+\sqrt{1+2A_k\sigma}\right)^2},
    \end{equation*}
    therefore for all $\sigma>0$, we have
    \begin{equation*}
    \begin{aligned}
        \left|\tau'(\sigma)\right| &\leq \frac{2A_k}{\left(1+\sqrt{1+2A_k\sigma}\right)^2}\\
        &\leq \frac{2A_k}{2A_k\sigma}\\
        &= \frac{1}{\sigma}.
    \end{aligned}
    \end{equation*}
    Therefore we have
    \begin{equation*}
        \|y_k(s)-y_k(t)\|\leq M_0\left|\tau'(\xi)\right|(s-t)\leq \frac{M_0}{t}(s-t).
    \end{equation*}
\end{proof}

\section{Alternative}
\begin{algorithm}[ht]
    \caption{Accelerated Trust-Region(ATR) Method}\label{alg.accelerated utr}
    \begin{algorithmic}[1]
        \STATE{\textbf{input:} $x_0=v_0\in \mathbb{R}^n$, $s_0=0$;}
        \FOR{$k=0, 1, \ldots, \infty$}
        \STATE{$y_k = \frac{k}{k+3}x_k+\frac{3}{k+3}v_k$}
        \STATE{$(d_k,\lambda_k) = \operatorname{UTR}(y_k,\frac{\sqrt{M}}{3}\|\nabla f(y_k)\|^{1/2},\frac{1}{3\sqrt{M}}\|\nabla f(y_k)\|^{1/2})$}
        \IF{$\lambda_k>0$}
        \STATE{$x_{k+1}=y_k+d_k$}
        \STATE{$s_k = s_{k-1}+\frac{(k+1)(k+2)}{2}\nabla f(y_k)$}
        \ELSE
        \STATE{$d_k^{(0)}=d_k$, $\lambda_k^{(0)} = \lambda_k$, $j=0$}
        \WHILE{$\lambda_k^{(j)}=0$}
        \STATE{$j=j+1$}
        \STATE{$(d_k^{(j)},\lambda_k^{(j)})=\operatorname{UTR}(y_k,\frac{\sqrt{M}}{3}\|\nabla f(y_k+d_k^{(j-1)})\|^{1/2},\frac{1}{3\sqrt{M}}\|\nabla f(y_k+d_k^{(j-1)})\|^{1/2})$}
        \ENDWHILE
        \IF{$\lambda_k^{(j)}\leq \frac{\sqrt{M}}{3}\|\nabla f(y_k+d_k^{(j-1)})\|^{1/2}$}
        % \STATE{$x_{k+1} = y_k+d_k^{(j)}$, $\mu_k = \lambda_k^{(j)}+\frac{\sqrt{M}}{3}\|\nabla f(y_k+d_k^{(j-1)})\|^{1/2}$}
        \STATE{$x_{k+1} = y_k+d_k^{(j)}$}
        \label{line.adaptive search output}
        \ELSE
        \STATE{$\mu_- = \frac{\sqrt{M}}{3}\|\nabla f(y_k+d_k^{(j-1)})\|^{1/2}$, $\mu_+ = \lambda_k^{(j)}+\frac{\sqrt{M}}{3}\|\nabla f(y_k+d_k^{(j-1)})\|^{1/2}$}
        % \STATE{$x_{k+1},\mu_k = \operatorname{BISECTION\_ATR}(y_k,\mu_-,\mu_+)$}
        \STATE{$x_{k+1} = \operatorname{BISECTION\_ATR}(y_k,\mu_-,\mu_+)$}
        \ENDIF
        \STATE{$s_k=s_{k-1}+\frac{(k+1)(k+2)}{2}\nabla f(x_{k+1})$}
        \ENDIF
        \STATE{$v_{k+1} = v_0-\sqrt{\frac{\ell}{3M\|s_k\|}}s_k$}
        \ENDFOR
    \end{algorithmic}
\end{algorithm}
\end{document}

